\SourcePage{33}
\SGAsection{5}{$R^qu_*$ 与拓扑斯态射的 Cartan--Leray 谱序列}
\textbf{5.0。}设 $u:(E,A)\to(E',A')$ 是\Term{ringed-topos}{赋环拓扑斯}的\Term{morphism}{态射}，$u_*:E\to E'$ 是\Term{direct-image}{正像函子}。仍用 $u_*$ 表示\Term{direct-image}{正像}\Term{functor}{函子}向\Term{module}{模}的延拓。\Term{module}{模}上的\Term{functor}{函子} $u_*$（为明确起见，取左\Term{module}{模}）是\Term{left-exact}{左正合}的。它的\Term{right-derived-functor}{右导出函子}记作 $R^qu_*$，$q\geq0$。
\begin{SGAresult}{命题}{5.1}
1）对每个 $A$-\Term{module}{模} $M$，\Term{sheaf}{层} $R^qu_*M$ 是\Term{presheaf}{预层} $X'\mapsto H^q(u^*X',M)$（$X'\in\operatorname{ob}E'$）的\Term{associated-sheaf}{伴随层}。

2）\Term{functor}{函子} $R^qu_*$ 的形成与\Term{restriction-scalars}{限制标量}交换。

3）\Term{functor}{函子} $R^qu_*$ 的形成与\Term{localization}{局部化}交换。确切地说，对 $E'$ 的每个\Term{object}{对象} $X'$，若以 $u/X':E/u^*X'\to E'/X'$ 表示由 $u$ \Term{localization}{局部化}得到的\Term{morphism}{态射}（第 IV 讲 8），则对每个 $A$-\Term{module}{模} $M$，都有\Term{canonical-isomorphism}{典范同构}
\begin{equation}R^q(u/X')_*(M|u^*X')\simeq R^qu_*(M)|X',\qquad q\geq0.\tag{5.1.1}\end{equation}
\end{SGAresult}
\SourcePage{34}
\begin{proof}
以 $\widehat u_*:\widehat E\to\widehat {E'}$ 表示 $\mathcal U$-\Term{presheaf}{预层}的\Term{direct-image}{正像函子}（$\widehat u_*M=M\circ u^*$）。因为 $u^*$ 与 $u_*$ 互为\Term{adjoint-functor}{伴随函子}，有\Term{isomorphism}{同构}
\[u_*\simeq a\widehat u_*^{\circ},\]
这里 $a$ 是 $E'$ 上的\Term{associated-sheaf}{伴随层函子}。由于\Term{functor}{函子} $a$ 与 $\widehat u_*$ 都是\Term{exact-functor}{正合的}，有
\[R^qu_*\simeq a\widehat u_*\mathscr H^q,\]
这就是 1）的另一种表述。由 1）与 3.5 得到 2）。根据\Term{morphism}{态射} $u/X'$ 的定义，$q=0$ 时有\Term{canonical-isomorphism}{典范同构} (5.1.1)。注意到\Term{localization}{局部化函子}（第 IV 讲 8）是\Term{exact}{正合}的，并把\Term{injective-object}{内射对象}变为\Term{injective-object}{内射对象}（4.11），便得到一般情形。
\end{proof}
\begin{SGAresult}{命题}{5.2}
设 $u:E\to E'$ 是\Term{topos}{拓扑斯}的\Term{morphism}{态射}，$S'$ 是 $E'$ 的一个\Term{generating-family}{生成族}。对所有\Term{functor}{函子} $H^0(u^*X',\cdot)$（$X'\in S'$）均\Term{acyclic}{无上同调}的\Term{sheaf}{层} $M$，对于\Term{functor}{函子} $u_*$ 也是\Term{acyclic}{无上同调}的。特别地，\Term{flasque-sheaf}{松层}对于 $u_*$ 是\Term{acyclic}{无上同调}的。
\end{SGAresult}
\begin{proof}由 5.1 的 1）得到。\end{proof}
\begin{SGAresult}{命题}{5.3}
设 $u:E\to E'$ 是\Term{topos}{拓扑斯}的\Term{morphism}{态射}，$M$ 是 $E$ 的一个\Term{abelian-group}{阿贝尔群}。有\Term{spectral-sequence}{谱序列}
\begin{equation}E_2^{p,q}=H^p(E',R^qu_*M)\Longrightarrow H^{p+q}(E,M).\tag{5.3.1}\end{equation}
更一般地，对 $E'$ 的每个\Term{object}{对象} $X'$，有\Term{spectral-sequence}{谱序列}
\begin{equation}E_2^{p,q}=H^p(X',R^qu_*M)\Longrightarrow H^{p+q}(u^*X',M).\tag{5.3.2}\end{equation}
\end{SGAresult}
\begin{proof}根据\Term{direct-image}{正像}与\Term{inverse-image}{逆像函子}的定义，有\Term{isomorphism}{同构} $H^0(X',u_*M)\simeq H^0(u^*X',M)$。\Term{functor}{函子} $u_*$ 把\Term{injective-object}{内射对象}变为\Term{flasque-sheaf}{松层}（4.6 与 4.9）。因此所述\Term{spectral-sequence}{谱序列}就是复合\Term{functor}{函子}的\Term{spectral-sequence}{谱序列}（0.3）。\end{proof}
\begin{SGAresult}{命题}{5.4}
设 $u:E\to E'$ 与 $v:E'\to E''$ 是两个\Term{topos}{拓扑斯}的\Term{morphism}{态射}。有\Term{spectral-sequence}{谱序列}
\[R^pv_*R^qu_*\Longrightarrow R^{p+q}(v\circ u)_*.\]
\end{SGAresult}
\begin{proof}有 $v_*u_*\simeq(vu)_*$，而\Term{functor}{函子} $u_*$ 把\Term{injective-object}{内射对象}变为\Term{flasque-sheaf}{松层}（4.9），因而变为对于 $v_*$ \Term{acyclic}{无上同调}的\Term{sheaf}{层}（5.2）。于是得到复合\Term{functor}{函子}的\Term{spectral-sequence}{谱序列}（0.3）。\end{proof}
\SourcePage{35}
\SGAsection{6}{局部 Ext 与带支集上同调}
\textbf{6.0。}设 $(E,A)$ 是\Term{ringed-topos}{赋环拓扑斯}，$M$ 是一个 $A$-\Term{module}{模}；为明确起见，取左\Term{module}{模}。在左 $A$-\Term{module}{模}\Term{category}{范畴}上、取值于\Term{abelian-group}{阿贝尔群}\Term{category}{范畴}的\Term{functor}{函子} $N\mapsto\mathscr Hom_A(M,N)$ 是\Term{left-exact}{左正合}的（第 IV 讲 12）。它的\Term{right-derived-functor}{右导出函子}记作
\begin{equation}\mathscr Ext_A^q(M,N).\tag{6.0.1}\end{equation}
特别地，有
\begin{equation}\mathscr Ext_A^0(M,N)=\mathscr Hom_A(M,N).\tag{6.0.2}\end{equation}
根据定义，有\Term{canonical-isomorphism}{典范同构}（2.1 与第 IV 讲 12）
\begin{equation}H^0(E,\mathscr Ext_A^0(M,N))=\operatorname{Ext}_A^0(E;M,N)=\operatorname{Hom}_A(M,N).\tag{6.0.3}\end{equation}
更一般地，对 $E$ 的每个\Term{object}{对象} $X$，有（2.2）
\begin{equation}H^0(X,\mathscr Ext_A^0(M,N))=\operatorname{Ext}_A^0(X;M,N)=\operatorname{Hom}_{A|X}(M|X,N|X).\tag{6.0.4}\end{equation}
\begin{SGAresult}{命题}{6.1}
1）\Term{functor}{函子} $\mathscr Ext_A^q$ 的形成与\Term{localization}{局部化}交换。确切地说，对 $E$ 的每个\Term{object}{对象} $X$，有\Term{functor}{函子}性的\Term{isomorphism}{同构}
\begin{equation}\mathscr Ext_A^q(M,N)|X\simeq\mathscr Ext_{A|X}^q(M|X,N|X).\tag{6.1.1}\end{equation}
2）\Term{sheaf}{层} $\mathscr Ext_A^q(M,N)$ \Term{isomorphism}{同构}于\Term{presheaf}{预层} $X\mapsto\operatorname{Ext}_A^q(X;M,N)$ 的\Term{associated-sheaf}{伴随层}。

3）存在\Term{spectral-sequence}{谱序列}
\begin{equation}\operatorname{Ext}_A^{p+q}(E;M,N)\Longleftarrow E_2^{p,q}=H^p(E,\mathscr Ext_A^q(F,G)).\tag{6.1.2}\end{equation}
更一般地，对 $E$ 的每个\Term{object}{对象} $X$，存在关于 $X$ 及变量 $M,N$ \Term{functor}{函子}性的\Term{spectral-sequence}{谱序列}
\begin{equation}\operatorname{Ext}_A^{p+q}(X;M,N)\Longleftarrow E_2^{p,q}=H^p(X,\mathscr Ext_A^q(M,N)).\tag{6.1.3}\end{equation}
\end{SGAresult}
\begin{proof}
根据定义，$q=0$ 时有\Term{isomorphism}{同构} (6.1.1)（第 IV 讲 12）。由于\Term{localization}{局部化}\Term{functor}{函子}是\Term{exact}{正合}的，并将\Term{injective-module}{内射模}变为\Term{injective-module}{内射模}（2.2），便得到一般情形。\Term{functor}{函子} $N\mapsto\mathscr Hom_A(M,N)=\mathscr Ext_A^0(M,N)$ 把\Term{injective-module}{内射模}变为\Term{flasque-sheaf}{松层}，因而变为对于 $H^0(X,\cdot)$ \Term{acyclic}{无上同调}的\Term{sheaf}{层}（4.10）。结合 (6.0.3) 与 (6.0.4)，\Term{spectral-sequence}{谱序列} (6.1.2)、(6.1.3) 就是复合\Term{functor}{函子}的\Term{spectral-sequence}{谱序列}。当 $X$ 在 $E$ 中变动时，\Term{spectral-sequence}{谱序列} (6.1.3) 给出\Term{presheaf}{预层}的\Term{spectral-sequence}{谱序列}，
\SourcePage{36}
取\Term{associated-sheaf}{伴随层}便得到\Term{sheaf}{层}的\Term{spectral-sequence}{谱序列}。由于 $p\ne0$ 时，\Term{presheaf}{预层} $X\mapsto H^p(X,\cdot)$ 的\Term{associated-sheaf}{伴随层}为零（3.1），这一\Term{spectral-sequence}{谱序列}\Term{degeneration}{退化}，并给出 2）中所述的\Term{isomorphism}{同构}。
\end{proof}
\begin{SGAresult}{命题}{6.2}
\Term{functor}{函子} $(M,N)\mapsto\mathscr Ext_A^q(M,N)$，$q\geq0$，在变量 $M$ 和变量 $N$ 上分别构成\Term{delta-functor}{$\delta$-函子}。明确地说，对每个\Term{exact-sequence}{正合列} $0\to N'\to N\to N''\to0$（相应地 $0\to M'\to M\to M''\to0$），有\Term{long-exact-sequence}{长正合列}
\begin{equation}\cdots\to\mathscr Ext_A^q(M,N')\to\mathscr Ext_A^q(M,N)\to\mathscr Ext_A^q(M,N'')\xrightarrow{\delta}\mathscr Ext_A^{q+1}(M,N')\to\cdots\tag{6.2.1}\end{equation}
（相应地
\begin{equation}\cdots\to\mathscr Ext_A^q(M'',N)\to\mathscr Ext_A^q(M,N)\to\mathscr Ext_A^q(M',N)\xrightarrow{\delta}\mathscr Ext_A^{q+1}(M'',N)\to\cdots).\tag{6.2.2}\end{equation}
\end{SGAresult}
\begin{proof}由 $\operatorname{Ext}^q$ 的一般性质，对 $E$ 的每个\Term{object}{对象} $X$，\Term{functor}{函子} $(M,N)\mapsto\operatorname{Ext}_A^q(X;M,N)$，$q\geq0$，在每个变量上都构成 $\delta$-\Term{functor}{函子}。令 $X$ 变动并取\Term{associated-sheaf}{伴随层}（6.1），即得结论。\end{proof}
\textbf{6.3。}设 $(E,A)$ 是\Term{ringed-topos}{赋环拓扑斯}，$M$ 是 $A$-\Term{module}{模}，$Z$ 是 $E$ 的\Term{closed-subtopos}{闭子拓扑斯}（第 IV 讲 9），$U$ 是它的补\Term{open-subtopos}{开子拓扑斯}。以 $H_Z^0(E,M)$ 表示 $M$ 的\Term{support}{支集}包含于 $Z$ 的\Term{section}{截面}所成的群（第 IV 讲 14），以 $\mathscr H_Z^0(M)$ 表示由 $M$ 的“\Term{support}{支集}在 $Z$ 中”的\Term{section}{截面}定义的 $M$ 的子\Term{sheaf}{层}（第 IV 讲 14）。\Term{functor}{函子} $H_Z^0(E,\cdot)$ 与 $\mathscr H_Z^0(\cdot)$ 是\Term{left-exact}{左正合}的（第 IV 讲 14）。它们的\Term{derived-functor}{导出函子}分别记作 $H_Z^q(E,\cdot)$ 与 $\mathscr H_Z^q(\cdot)$，称为 $M$ 的\Term{cohomology-supports}{支集在 $Z$ 中的上同调群}（相应地，\Term{cohomology}{上同调}\Term{sheaf}{层}）。\footnote{普通\Term{topological-space}{拓扑空间}的情形，可比较 SGA 2 第 I 讲，以及下文第 80 页引用的 Hartshorne 的讲义。}
有\Term{canonical-isomorphism}{典范同构}（第 IV 讲 14）
\begin{align}H_Z^0(E,M)&\simeq\operatorname{Hom}_A(A_Z,M)\simeq\operatorname{Ext}_A^0(E;A_Z,M),\tag{6.3.1}\\
\mathscr H_Z^0(M)&\simeq\mathscr Hom_A(A_Z,M)\simeq\mathscr Ext_A^0(A_Z,M),\tag{6.3.2}\end{align}
从而对每个 $q\geq0$，有\Term{isomorphism}{同构}
\begin{align}H_Z^q(E,M)&\simeq\operatorname{Ext}_A^q(E;A_Z,M),\tag{6.3.3}\\
\mathscr H_Z^q(M)&\simeq\mathscr Ext_A^q(A_Z,M).\tag{6.3.4}\end{align}
注意 $A_Z$ 是\Term{bimodule}{双模}，所以\Term{sheaf}{层} $\mathscr Ext_A^q(A_Z,M)\simeq\mathscr H_Z^q(M)$ \Term{canonical}{典范}地具有 $A$-\Term{module}{模}结构。

对 $E$ 的每个\Term{object}{对象} $X$，以 $Z/X$ 表示 $E/X$ 中由 $Z$ \Term{localization}{局部化}得到的\Term{closed-subtopos}{闭子拓扑斯}（即 $U\times X$ 的补）。根据定义，有\Term{canonical-isomorphism}{典范同构}
\begin{equation}H^0(X,\mathscr H_Z^0(M))\simeq H_{Z/X}^0(E/X,M|X).\tag{6.3.5}\end{equation}
\SourcePage{37}
置
\begin{equation}H_Z^q(X,M)\simeq H_{Z/X}^q(E/X,M|X).\tag{6.3.6}\end{equation}
结合 (6.3.2)，有\Term{canonical-isomorphism}{典范同构}
\begin{equation}H_Z^q(X,M)\simeq\operatorname{Ext}_A^q(X;A_Z,M).\tag{6.3.7}\end{equation}
\begin{SGAresult}{命题}{6.4}
1）\Term{functor}{函子} $\mathscr H_Z^q$ 的形成与\Term{localization}{局部化}交换。确切地说，对 $E$ 的每个\Term{object}{对象} $X$ 及每个 $A$-\Term{module}{模} $M$，有\Term{canonical-isomorphism}{典范同构}
\begin{equation}\mathscr H_Z^0(M)|X\simeq\mathscr H_{Z/X}^q(M|X).\tag{6.4.1}\end{equation}
2）\Term{sheaf}{层} $\mathscr H_Z^q(M)$ 是\Term{presheaf}{预层} $X\mapsto H_Z^q(X,M)$ 的\Term{associated-sheaf}{伴随层}。

3）存在\Term{spectral-sequence}{谱序列}
\begin{equation}H_Z^{p+q}(E,M)\Longleftarrow H^p(E,\mathscr H_Z^q(M)).\tag{6.4.2}\end{equation}
更一般地，对 $E$ 的每个\Term{object}{对象} $X$，存在\Term{spectral-sequence}{谱序列}
\begin{equation}H_Z^{p+q}(X,M)\Longleftarrow H^p(X,\mathscr H_Z^q(M)).\tag{6.4.3}\end{equation}
\end{SGAresult}
\begin{proof}只需利用 6.3 的对应关系翻译命题 6.1。\end{proof}
\begin{SGAresult}{命题}{6.5}
沿用 6.3 的记号，以 $j:U\to E$ 表示\Term{canonical}{典范}\Term{morphism}{态射}。对每个 $A$-\Term{module}{模} $M$，存在\Term{sheaf}{层}的\Term{exact-sequence}{正合列}
\begin{equation}0\to\mathscr H_Z^0(M)\to M\to j_*(M|U)\to\mathscr H_Z^1(M)\to0,\tag{6.5.1}\end{equation}
而对 $q\geq2$，有\Term{isomorphism}{同构}
\begin{equation}\mathscr H_Z^q(M)\simeq\mathscr Ext_A^{q-1}(A_U,M)\simeq R^{q-1}j_*(M|U).\tag{6.5.2}\end{equation}
此外，有\Term{long-exact-sequence}{长正合列}
\begin{equation}\cdots\to H_Z^q(E,M)\to H^q(E,M)\to H^q(U,M)\to H_Z^{q+1}(E,M)\to\cdots,\tag{6.5.3}\end{equation}
更一般地，对 $E$ 的每个\Term{object}{对象} $X$，有\Term{long-exact-sequence}{长正合列}
\begin{equation}\cdots\to H_Z^q(X,M)\to H^q(X,M)\to H^q(X\times U,M)\to H_Z^{q+1}(X,M)\to\cdots.\tag{6.5.4}\end{equation}
\end{SGAresult}
\footnote{此\Term{exact-sequence}{正合列}明确了整体\Term{cohomology}{上同调}\Term{invariants}{不变量} $H_Z^q(E,N)$ 作为“$E$ 模去\Term{open-subset}{开集} $U$、以 $M$ 为系数的\Term{cohomology}{上同调}群”的\Term{action}{作用}。}
\begin{proof}根据 $A_Z$ 的定义，有\Term{exact-sequence}{正合列}（第 IV 讲 14）
\begin{equation}0\to A_U\to A\to A_Z\to0.\tag{6.5.5}\end{equation}
$q>0$ 时，\Term{functor}{函子} $\mathscr Ext_A^q(A,\cdot)$ 为零。有 $\mathscr Hom_A(A_U,M)\simeq j_*(M|U)$
\SourcePage{38}
（第 IV 讲 14），从而（2.2）$\mathscr Ext_A^q(A_U,M)\simeq R^qj_*(M|U)$。最后 $\mathscr Ext_A^q(A_Z,M)\simeq\mathscr H_Z^q(M)$（6.3.4）。\Term{long-exact-sequence}{长正合列} (6.2.2) 在此情形给出 (6.5.1) 与 (6.5.2)。\Term{exact-sequence}{正合列} (6.5.3)、(6.5.4) 由 6.3 的对应关系以及 $\delta$-\Term{functor}{函子} $\operatorname{Ext}_A^q(X;\cdot,M)$，$q\geq0$，关于 (6.5.5) 的\Term{long-exact-sequence}{长正合列}得到。\end{proof}
\begin{SGAresult}{命题}{6.6}
1）\Term{flasque-sheaf}{松层}对于\Term{functor}{函子} $\mathscr H_Z^0$ 和 $H_Z^0(X,\cdot)$ 是\Term{acyclic}{无上同调}的。

2）\Term{functor}{函子} $\mathscr H_Z^0$ 和 $H_Z^0(X,\cdot)$ 与\Term{restriction-scalars}{限制标量}交换。
\end{SGAresult}
\begin{proof}
设 $M$ 为\Term{flasque-sheaf}{松层}。\Term{exact-sequence}{正合列} (6.5.4) 给出：对一切 $X$ 与一切 $q\geq2$，有 $H_Z^q(X,M)=0$，并有\Term{exact-sequence}{正合列}
\[0\to H_Z^0(X,M)\to H^0(X,M)\to H^0(X\times U,M)\to H_Z^1(X,M)\to0.\]
但因\Term{sheaf}{层} $M$ 是松的，\Term{morphism}{态射} $H^0(X,M)\to H^0(X\times U,M)$ 是\Term{surjectivity}{满射}（4）。所以 $H_Z^1(X,M)=0$，且 $M$ 对于 $H_Z^0(X;M)$ 是\Term{acyclic}{无上同调}的。取\Term{associated-sheaf}{伴随层}即知，$M$ 对于\Term{functor}{函子} $\mathscr H_Z^0$ 是\Term{acyclic}{无上同调}的（6.4）。\Term{functor}{函子} $\mathscr H_Z^0$ 与 $H_Z^0(X,\cdot)$ 显然与\Term{restriction-scalars}{限制标量}交换；由于\Term{flasque-sheaf}{松层}对这两个\Term{functor}{函子}均\Term{acyclic}{无上同调}，性质 2）随之成立。
\end{proof}
\begin{SGAresult}{命题}{6.7}
设 $(E,A)$ 是\Term{ringed-topos}{赋环拓扑斯}，$Z$ 是 $E$ 的\Term{closed-subtopos}{闭子拓扑斯}，$U$ 是补\Term{open-subtopos}{开子拓扑斯}，$M,N$ 是两个 $A$-\Term{module}{模}。存在关于 $M,N$ \Term{functor}{函子}性的、与\Term{closed-subtopos}{闭子拓扑斯}的改变相容的\Term{isomorphism}{同构}
\[\mathscr H_Z^0(\mathscr Hom_A(M,N))\simeq\mathscr Hom_A(M,\mathscr H_Z^0(N))\simeq\mathscr Hom_A(M\otimes_A A_Z,N).\]
\end{SGAresult}
\begin{proof}由第 IV 讲 14 结合 (6.3.2) 得到。\end{proof}
\textbf{6.8。}置
\begin{equation}\mathscr Hom_{A,Z}(M,N)=\mathscr H_Z^0(\mathscr Hom_A(M,N)).\tag{6.8.1}\end{equation}
\Term{functor}{函子} $N\mapsto\mathscr Hom_{A,Z}(M,N)$ 是\Term{left-exact}{左正合}的。它的\Term{derived-functor}{导出函子}记作 $\mathscr Ext_{A,Z}^q(M,N)$：这些就是\Term{ext-supports}{支集在 $Z$ 中的 Ext 层}。因此有
\begin{equation}\mathscr Ext_{A,Z}^0(M,N)=\mathscr Hom_{A,Z}(M,N).\tag{6.8.2}\end{equation}
置 $M_Z=M\otimes_A A_Z$。由第 IV 讲 14，$M_Z$ 是 $M$ 在 $Z$ 上的\Term{inverse-image}{逆像}在 $E$ 上的\Term{direct-image}{正像}。因而由 6.7 有\Term{isomorphism}{同构}
\begin{equation}\mathscr Ext_{A,Z}^q(M,N)\simeq\mathscr Ext_A^q(M_Z,N).\tag{6.8.3}\end{equation}
现在转向整体\Term{invariants}{不变量}。置
\begin{equation}\operatorname{Hom}_{A,Z}(M,N)=\operatorname{Ext}_{A,Z}^0(E;M,N)=H_Z^0(E,\mathscr Hom_A(M,N)).\tag{6.8.4}\end{equation}
群 $\operatorname{Hom}_{A,Z}(M,N)$ 是从 $M$ 到 $N$ 的\Term{morphism}{态射}群中\Term{support}{支集}在 $Z$ 内（第 IV 讲 14），即在 $U$ 上为零的\Term{morphism}{态射}所成的\Term{subgroup}{子群}。更一般地，对 $E$ 的每个\Term{object}{对象} $X$，置
\begin{equation}\operatorname{Ext}_{A,Z}^0(X;M,N)=H_Z^0(X,\mathscr Hom_A(M,N)).\tag{6.8.5}\end{equation}
\SourcePage{39}
\Term{functor}{函子} $N\mapsto\operatorname{Ext}_{A,Z}^0(X;M,N)$ 是\Term{left-exact}{左正合}的。\Term{derived-functor}{导出函子}记作 $\operatorname{Ext}_{A,Z}^q(X;M,N)$，称为\Term{ext-supports}{支集在 $Z$ 中的 Ext 群}。定义 (6.3.6)、(6.8.1)、(6.8.5) 及\Term{isomorphism}{同构} 6.7 给出\Term{isomorphism}{同构}
\begin{equation}\operatorname{Ext}_{A,Z}^0(X;M,N)\simeq
\begin{cases}H^0(X,\mathscr Hom_{A,Z}(M,N)),\\
\operatorname{Ext}^0(X;M,\mathscr H_Z^0(N)),\\
\operatorname{Ext}^0(X;M_Z,N).
\end{cases}\tag{6.8.6}\end{equation}
(6.8.6) 的最后一个\Term{isomorphism}{同构}给出\Term{isomorphism}{同构}
\begin{equation}\operatorname{Ext}_{A,Z}^q(X;M,N)\simeq\operatorname{Ext}_A^q(X;M_Z,N).\tag{6.8.7}\end{equation}
\begin{SGAresult}{命题}{6.9}
1）存在两个关于 $M,N$ \Term{functor}{函子}性的、与\Term{closed-subtopos}{闭子拓扑斯}的改变相容的\Term{spectral-sequence}{谱序列}
\begin{equation}\mathscr Ext_{A,Z}^{p+q}(M,N)\Longleftarrow
\begin{cases}E_2^{p,q}=\mathscr H_Z^p(\mathscr Ext_A^q(M,N)),\\
{}'E_2^{p,q}=\mathscr Ext_A^p(M,\mathscr H_Z^q(N)).\end{cases}\tag{6.9.1}\end{equation}
2）存在三个关于 $X,M,N$ \Term{functor}{函子}性的、与\Term{closed-subtopos}{闭子拓扑斯}的改变相容的\Term{spectral-sequence}{谱序列}
\begin{equation}\operatorname{Ext}_{A,Z}^{p+q}(M,N)\Longleftarrow
\begin{cases}E_2^{p,q}=H_Z^p(X,\mathscr Ext_A^q(M,N)),\\
{}'E_2^{p,q}=H^p(X,\mathscr Ext_{A,Z}^q(M,N)),\\
{}''E_2^{p,q}=\operatorname{Ext}_A^p(X;M,\mathscr H_Z^q(N)).\end{cases}\tag{6.9.2}\end{equation}
3）\Term{sheaf}{层} $\mathscr Ext_{A,Z}^q(M,N)$ \Term{canonical-isomorphism}{典范同构}于\Term{presheaf}{预层} $X\mapsto\operatorname{Ext}_{A,Z}^q(X;M,N)$ 的\Term{associated-sheaf}{伴随层}。
\end{SGAresult}
\begin{proof}
当 $N$ \Term{injective}{内射}时，\Term{sheaf}{层} $\mathscr Hom_A(M,N)$ 是松的（4.10），所以对于 $\mathscr H_Z^0$ \Term{acyclic}{无上同调}（6.6）。(6.9.1) 的第一个\Term{spectral-sequence}{谱序列}就是复合\Term{functor}{函子}的\Term{spectral-sequence}{谱序列}（6.8.1）。同样，当 $N$ \Term{injective}{内射}时，$\mathscr H_Z^0(N)$ \Term{injective}{内射}（4.11），而 (6.9.1) 的第二个\Term{spectral-sequence}{谱序列}是从 6.7 得到的复合\Term{functor}{函子}\Term{spectral-sequence}{谱序列}。(6.9.2) 的\Term{spectral-sequence}{谱序列}是由定义 (6.8.5) 与 (6.8.6) 的头两个\Term{isomorphism}{同构}得到的复合\Term{functor}{函子}\Term{spectral-sequence}{谱序列}。最后，令 (6.9.2) 的第二个\Term{spectral-sequence}{谱序列}中的 $X$ 变动并取\Term{associated-sheaf}{伴随层}，得到\Term{sheaf}{层}的\Term{spectral-sequence}{谱序列}；它由 3.1 \Term{degeneration}{退化}，给出 3）中的\Term{isomorphism}{同构}。
\end{proof}
\begin{SGAresult}{命题}{6.10}
\Term{functor}{函子} $(M,N)\mapsto\mathscr Ext_{A,Z}^q(M,N)$，$q\geq0$，以及 $(M,N)\mapsto\operatorname{Ext}_{A,Z}^q(X;M,N)$，$q\geq0$，在变量 $M,N$ 上分别都是 $\delta$-\Term{functor}{函子}。
\SourcePage{40}
以 $M_U$ 表示\Term{sheaf}{层} $M\otimes_A A_U$（参见第 IV 讲 11），有\Term{long-exact-sequence}{长正合列}
\begin{equation}\cdots\xrightarrow{\delta}\mathscr Ext_{A,Z}^q(M,N)\to\mathscr Ext_A^q(M,N)\to\mathscr Ext_A^q(M_U,N)\xrightarrow{\delta}\mathscr Ext_{A,Z}^{q+1}(M,N)\to\cdots\tag{6.10.1}\end{equation}
以及 Ext 群的类似\Term{exact-sequence}{正合列}。
\end{SGAresult}
\begin{proof}
\Term{functor}{函子} $\mathscr Ext_A^q(\cdot,\cdot)$ 与 $\operatorname{Ext}_A^q(X;\cdot,\cdot)$ 关于每个变量分别构成 $\delta$-\Term{functor}{函子}（6.2），而由于 $A_Z$ \Term{flat}{平坦}（1.3.3），\Term{functor}{函子} $M\mapsto M_Z$ 是\Term{exact}{正合}的。因此第一个断言由\Term{isomorphism}{同构} (6.8.3) 与 (6.8.7) 得到。结合 (6.8.3)、(6.8.7)，\Term{exact-sequence}{正合列} (6.10.1) 及 Ext 群的类似\Term{exact-sequence}{正合列}，就是 $\delta$-\Term{functor}{函子}
\[\mathscr Ext_A^q(\cdot,N),\ q\geq0\quad\text{（相应地 }\operatorname{Ext}_A^q(X;\cdot,N),\ q\geq0\text{）}\]
关于\Term{exact-sequence}{正合列} $0\to M_U\to M\to M_Z\to0$ 的\Term{long-exact-sequence}{长正合列}。
\end{proof}
\textbf{6.11。}简略说明如何把这些结果推广到\Term{family-supports}{支集族}的情形。设 $E$ 是\Term{topos}{拓扑斯}。对 $E$ 的每个\Term{object}{对象} $X$，以 $\operatorname{Fer}(X)$ 表示\Term{topos}{拓扑斯} $E/X$ 的\Term{closed-subtopos}{闭子拓扑斯}所成的\Term{set}{集合}。取补集给出此\Term{set}{集合}与 $E/X$ 的\Term{open-subtopos}{开子拓扑斯}所成的\Term{set}{集合}，即 $X$ 的\Term{subobject}{子对象}\Term{set}{集合}（第 IV 讲 8）之间的一一对应。容易立即验证，\Term{presheaf}{预层} $\operatorname{Fer}:X\mapsto\operatorname{Fer}(X)$ 实际上是 $\mathcal U$-\Term{sheaf}{层}，因此\Term{representable-functor}{可表示}。换言之，存在 $E$ 的一个\Term{object}{对象}，仍记作 $\operatorname{Fer}$，以及 $E/\operatorname{Fer}$ 的一个\Term{closed-subtopos}{闭子拓扑斯} $Z_{\operatorname{Fer}}$，使得对每个 $X$，$\operatorname{Fer}(X)$ 的每个元素 $Z$ 都可由 $Z_{\operatorname{Fer}}$ 沿一个由 $Z$ 唯一确定的\Term{morphism}{态射} $u_Z:X\to\operatorname{Fer}$ 作\Term{base-change}{基变换}得到。
\begin{SGAresult}{定义}{6.12}
$E$ 的\Term{family-supports}{支集族}，指 $E$ 的\Term{closed-subtopos}{闭子拓扑斯}\Term{set}{集合}的一个\Term{subset}{子集} $\Phi$，满足：

（S1）$\Phi$ 的\Term{finite}{有限}个元素的并仍属于 $\Phi$；

（S2）$E$ 中包含于 $\Phi$ 的某个元素的每个\Term{closed-subtopos}{闭子拓扑斯}，都属于 $\Phi$。
\end{SGAresult}
\textbf{6.12.1。}设 $E$ 是\Term{topos}{拓扑斯}，$\Phi$ 是 $E$ 的\Term{family-supports}{支集族}，$X$ 是 $E$ 的\Term{object}{对象}。以 $\Phi(X)$ 表示 $E/X$ 的最小\Term{family-supports}{支集族}，它包含由 $\Phi$ 的\Term{closed-subtopos}{闭子拓扑斯}沿\Term{base-change}{基变换} $X\to e$（$e$ 是 $E$ 的\Term{final-object}{终对象}）得到的 $E/X$ 的\Term{closed-subtopos}{闭子拓扑斯}。\Term{functor}{函子} $X\mapsto\Phi(X)$ 是\Term{sheaf}{层} $\operatorname{Fer}$ 的子\Term{presheaf}{预层}，因而是\Term{separated}{分离}的。

\textbf{6.12.2。}若 $E$ 的\Term{family-supports}{支集族} $\Phi$ 具有下列性质，就说它具有\Term{local-character}{局部性}：

（CL）对每个\Term{epimorphic-family}{满射族} $(X_i\to e)$，$i\in I$（其中 $e$ 是 $E$ 的\Term{final-object}{终对象}），\Term{set}{集合}列
\[\Phi\longrightarrow\prod_i\Phi(X_i)\rightrightarrows\prod_{i,j}\Phi(X_i\times X_j)\]
是\Term{exact}{正合}的。

设 $\widetilde\Phi$ 是\Term{presheaf}{预层} $X\mapsto\Phi(X)$ 的\Term{associated-sheaf}{伴随层}（6.12.1）。条件（CL）\Term{equivalence}{等价}于\Term{canonical}{典范}\Term{morphism}{态射} $\Phi\to\widetilde\Phi(e)$ 是\Term{bijection}{双射}（参见第 II 讲中\Term{separated}{分离}\Term{presheaf}{预层}的\Term{associated-sheaf}{伴随层}构造）。
\SourcePage{41}
因此验证（CL）时，可以只考虑\Term{epimorphic-family}{满射族} $(X_i\to e)$，$i\in I$，的一个终族。
\begin{SGAresult}{例}{6.12.3}
1）设 $Z$ 是 $E$ 的\Term{closed-subtopos}{闭子拓扑斯}。$E$ 中包含于 $Z$ 的\Term{closed-subtopos}{闭子拓扑斯}所成的\Term{set}{集合}是 $E$ 的\Term{family-supports}{支集族}，并具有\Term{local-character}{局部性}。

2）设 $T$ 为\Term{topological-space}{拓扑空间}，$\Phi$ 是 $T$ 的一个\Term{paracompactifying-supports}{仿紧化支集族} [7]。一般而言，$\Phi$ 不具有\Term{local-character}{局部性}。

3）设 $T$ 为\Term{topological-space}{拓扑空间}，$p$ 为整数。$T$ 中 Krull \Term{codimension}{余维} $\geq p$ 的闭集所成的族 $\Phi^p$ 是 $T$ 的\Term{family-supports}{支集族}，并具有\Term{local-character}{局部性}。

4）设\Term{topos}{拓扑斯} $E$ 具有如下性质：每个\Term{epimorphic-family}{满射族} $(X_i\to e)$，$i\in I$，都被一个\Term{finite}{有限}\Term{epimorphic-family}{满射族}\Term{refinement}{加细}。则 $E$ 的每个\Term{family-supports}{支集族}都具有\Term{local-character}{局部性}。事实上，由例 1），只需证明具有\Term{local-character}{局部性}的\Term{family-supports}{支集族} $\Phi_\lambda$ 的\Term{filtered-colimit}{滤过归纳极限}仍具有\Term{local-character}{局部性}；这立即由对\Term{set}{集合}列
\[\Phi_\lambda\longrightarrow\prod_i\Phi_\lambda(X_i)\rightrightarrows\prod_{i,j}\Phi_\lambda(X_i\times X_j)\]
取\Term{inductive-limit}{归纳极限}得到，其中 $(X_i\to e)$，$i\in I$，是\Term{finite}{有限}\Term{epimorphic-family}{满射族}。确实，由 6.12.2 和关于 $E$ 的假设，验证（CL）时可以只考虑\Term{finite}{有限}\Term{epimorphic-family}{满射族}，而\Term{filtered-colimit}{滤过归纳极限}与\Term{finite}{有限}\Term{projective-limit}{投射极限}交换（第 I 讲 2），所以取\Term{inductive-limit}{归纳极限}保留这些\Term{set}{集合}列的\Term{exactness}{正合性}。

5）设 $E$ 是\Term{coherent-topos}{凝聚拓扑斯}（第 VI 讲 2.3）。则由上文，$E$ 的每个\Term{family-supports}{支集族}都具有\Term{local-character}{局部性}。
\end{SGAresult}
\textbf{6.13。}设 $(E,A)$ 是\Term{ringed-topos}{赋环拓扑斯}，$\Phi$ 是 $E$ 的\Term{family-supports}{支集族}，$N$ 是 $A$-\Term{module}{模}（为明确起见，取左\Term{module}{模}），$X$ 是 $E$ 的\Term{object}{对象}。置
\begin{align}\mathscr H_\Phi^0(N)&=\varinjlim_{Z\in\Phi}\mathscr H_Z^0(N),\tag{6.13.1}\\
H_\Phi^0(E,N)&=\varinjlim_{Z\in\Phi}H_Z^0(E,N).\tag{6.13.2}\end{align}
设 $M$ 是左 $A$-\Term{module}{模}，置
\begin{align}\mathscr Ext_{A,\Phi}^0(M,N)=\mathscr Hom_{A,\Phi}(M,N)&=\varinjlim_{Z\in\Phi}\mathscr Ext_{A,Z}^0(M,N),\tag{6.13.3}\\
\operatorname{Ext}_{A,\Phi}^0(X;M,N)=\operatorname{Hom}_{A|X,\Phi(X)}(M|X,N|X)&=\varinjlim_{Z\in\Phi}\operatorname{Ext}_{A,Z}^0(X;M,N).\tag{6.13.4}\end{align}
有\Term{canonical-isomorphism}{典范同构}
\begin{equation}\mathscr Ext_{A,\Phi}^0(M,N)\simeq\mathscr H_\Phi^0(\mathscr Ext_A^0(M,N)),\tag{6.13.5}\end{equation}
\SourcePage{42}
\begin{equation}\operatorname{Ext}_{A,\Phi}^0(X;M,N)\simeq H^0((X,\mathscr Ext_A^0(M,N)).\tag{6.13.6}\end{equation}
当 $\Phi$ 是包含于一个\Term{closed-subtopos}{闭子拓扑斯} $Z$ 的\Term{closed-subtopos}{闭子拓扑斯}所成的族时，有\Term{isomorphism}{同构}
\begin{equation}\begin{cases}\mathscr H_\Phi^0\simeq\mathscr H_Z^0,\\ H_\Phi^0\simeq H_Z^0,\\
\mathscr Ext_{A,\Phi}^0(\cdot,\cdot)\simeq\mathscr Ext_{A,Z}^0(\cdot,\cdot),\\
\operatorname{Ext}_{A,\Phi}^0(X;\cdot,\cdot)\simeq\operatorname{Ext}_{A,Z}^0(X;\cdot,\cdot).
\end{cases}\tag{6.13.7}\end{equation}
定义上述\Term{functor}{函子}的\Term{inductive-limit}{归纳极限}是\Term{filtered-category}{滤过的}。因此这些\Term{functor}{函子}都是\Term{left-exact}{左正合}的。它们的\Term{right-derived-functor}{右导出函子}记作
\begin{equation}\mathscr H_\Phi^q,\ H_\Phi^q,\ \mathscr Ext_{A,\Phi}^q(\cdot,\cdot),\ \operatorname{Ext}_{A,\Phi}^q(X;\cdot,\cdot).\tag{6.13.8}\end{equation}
这些\Term{functor}{函子}同样可通过令 $Z$ 遍历 $\Phi$，对 $\mathscr H_Z^q,H_Z^q$ 等\Term{functor}{函子}取\Term{inductive-limit}{归纳极限}来计算。

由\Term{isomorphism}{同构} (6.13.5)、(6.13.6)，对 (6.9.1) 的第一个\Term{spectral-sequence}{谱序列}及 (6.9.2) 的第一个\Term{spectral-sequence}{谱序列}取\Term{inductive-limit}{归纳极限}，得到两个\Term{spectral-sequence}{谱序列}：
\begin{align}\mathscr Ext_{A,\Phi}^{p+q}(M,N)&\Longleftarrow E_2^{p,q}=\mathscr H_\Phi^p(\mathscr Ext_A^q(M,N)),\tag{6.13.9}\\
\operatorname{Ext}_{A,\Phi}^{p+q}(X;M,N)&\Longleftarrow E_2^{p,q}=H_\Phi^p(X,\mathscr Ext_A^q(M,N)).\tag{6.13.10}\end{align}
\textbf{6.14。}沿用 6.13 的记号，无须对 $\Phi$ 加任何假设，就有\Term{functor}{函子}性\Term{morphism}{态射}
\begin{equation}\theta:H_\Phi^0(N)\longrightarrow H^0(E,\mathscr H_\Phi^0(N)).\tag{6.14.1}\end{equation}
此\Term{morphism}{态射}总是\Term{injectivity}{单射}，但一般不是\Term{isomorphism}{同构}。不过，若 $\Phi$ 具有\Term{local-character}{局部性}，容易立即验证\Term{morphism}{态射} $\theta$ 是\Term{isomorphism}{同构}。同样，若 $\Phi$ 具有\Term{local-character}{局部性}，就有\Term{isomorphism}{同构}
\begin{equation}\theta':\operatorname{Ext}_{A,\Phi}^0(E;M,N)\simeq H^0(E,\mathscr Ext_{A,\Phi}^0(M,N)).\tag{6.14.2}\end{equation}
\begin{SGAresult}{命题}{6.15}
设 $(E,A)$ 是\Term{ringed-topos}{赋环拓扑斯}，$E$ \Term{coherent}{凝聚}（第 VI 讲），$\Phi$ 是 $E$ 的\Term{family-supports}{支集族}，$M,N$ 是两个 $A$-\Term{module}{模}。存在两个\Term{spectral-sequence}{谱序列}：
\begin{align}H_\Phi^{p+q}(E,N)&\Longleftarrow E_2^{p,q}=H_\Phi^p(E,\mathscr H_\Phi^q(N)),\tag{6.15.1}\\
\operatorname{Ext}_{A,\Phi}^{p+q}(E;M,N)&\Longleftarrow E_2^{p,q}=H^p(E,\mathscr Ext_{A,\Phi}^q(M,N)).\tag{6.15.2}\end{align}
\end{SGAresult}
\begin{proof}令\Term{closed-subtopos}{闭子拓扑斯} $Z$ 遍历 $\Phi$，对\Term{spectral-sequence}{谱序列} (6.4.2) 及 (6.9.2) 的第二个\Term{spectral-sequence}{谱序列}取\Term{inductive-limit}{归纳极限}，便得到这些\Term{spectral-sequence}{谱序列}；这里利用了\Term{coherent-topos}{凝聚拓扑斯}的\Term{cohomology}{上同调}与\Term{sheaf}{层}的\Term{inductive-limit}{归纳极限}交换这一事实（第 IV 讲 5）。\end{proof}
\SourcePage{43}
\textbf{6.16。}不加证明地指出：假设\Term{topos}{拓扑斯} $E$ \Term{coherent}{凝聚}且\Term{module}{模} $M$ 是\Term{perfect-module}{完美模} [13]，就能把 (6.9.1) 的第二个\Term{spectral-sequence}{谱序列}及 (6.9.2) 的第三个\Term{spectral-sequence}{谱序列}（取 $X$ 为 $E$ 的\Term{final-object}{终对象}）推广到\Term{family-supports}{支集族}的情形。

最后，还可以通过引入\Term{family-supports}{支集族}的\Term{presheaf}{预层}、\Term{family-supports}{支集族}的\Term{sheaf}{层}以及相应的\Term{cohomology}{上同调}群与\Term{cohomology}{上同调}\Term{sheaf}{层}，推广\Term{family-supports}{支集族}的概念。\footnote{这些概念以赋格及其变奏的形式展开，见 [12] 第 IV 章第 1 节（Lecture Notes 20，Springer）。}
\SGAsection{7}{附录：Čech 上同调}
本节发展 P. Cartier 的一个想法，说明在任意\Term{topos}{拓扑斯}中如何利用\Term{covering}{覆盖}计算\Term{sheaf}{层}的\Term{cohomology}{上同调}。仿\Term{compact-space}{紧空间}的经典理论见 [7]；适用于某些\Term{topos}{拓扑斯}，特别是\Term{etale}{平展}\Term{topos}{拓扑斯}的另一种方法，见 [14]。

\textbf{7.1。\Term{skeleton}{骨架}与\Term{coskeleton}{余骨架}。}

\textbf{7.1.0。}设 $\boldsymbol\Delta$ 是\Term{standard-simplex}{标准单形}\Term{category}{范畴}（它的\Term{object}{对象}是\Term{ordered-set}{有序集} $\Delta_n=[0,\ldots,n]$，\Term{morphism}{态射}是递增映射）。设 $E$ 是 $\mathcal U$-\Term{topos}{拓扑斯}，$\boldsymbol\Delta E$ 是 $\boldsymbol\Delta$ 上取值于 $E$ 的\Term{presheaf}{预层}\Term{category}{范畴}，即 $E$ 的\Term{semisimplicial-object}{半单纯对象}\Term{category}{范畴}。以 $\boldsymbol\Delta[n]$ 表示 $\boldsymbol\Delta$ 中由\Term{object}{对象} $\Delta_p$，$p\leq n$，定义的\Term{full-subcategory}{满子范畴}，以 $i_n:\boldsymbol\Delta[n]\to\boldsymbol\Delta$ 表示包含\Term{functor}{函子}，以 $\boldsymbol\Delta E[n]$ 表示 $\boldsymbol\Delta[n]$ 上取值于 $E$ 的\Term{presheaf}{预层}\Term{category}{范畴}，即 $E$ 的截断到 $n$ 阶的\Term{semisimplicial-object}{半单纯对象}\Term{category}{范畴}。由第 I 讲 5.1，有三个\Term{adjoint-functor}{伴随函子}构成的列（第 I 讲 5.3）
\[\begin{tikzcd}[column sep=large]
\boldsymbol\Delta E[n]\arrow[r,bend left=32,"i_{n!}"]\arrow[r,bend right=32,"i_{n*}"']&\boldsymbol\Delta E\arrow[l,"i_n^*" description]
\end{tikzcd}\]
其中 $i_n^*$ 是向\Term{category}{范畴} $\boldsymbol\Delta[n]$ 的\Term{restriction}{限制}\Term{functor}{函子}，而 $i_{n*}$ 与 $i_{n!}$ 分别是它的\Term{right-adjoint}{右伴随}与\Term{left-adjoint}{左伴随}。注意 $\boldsymbol\Delta E$ 与 $\boldsymbol\Delta E[n]$ 都是 $\mathcal U$-\Term{topos}{拓扑斯}（第 IV 讲 1.2），且 $(i_n^*,i_{n*})$ 是\Term{topos}{拓扑斯}\Term{morphism}{态射}（第 IV 讲 3.1），它给出 $\boldsymbol\Delta E[n]$ 到 $\boldsymbol\Delta E$ 的\Term{embedding}{嵌入}（第 I 讲 5.6 及第 IV 讲 9.1.1）。
\begin{SGAresult}{定义}{7.1.1}
$\boldsymbol\Delta E$ 到自身的\Term{functor}{函子} $i_{n!}i_n^*$（相应地 $i_{n*}i_n^*$）记作 $\operatorname{sk}_n$（相应地 $\operatorname{cosk}_n$），称为 $n$ 阶\Term{skeleton}{骨架函子}（相应地，$n$ 阶\Term{coskeleton}{余骨架函子}）。
\end{SGAresult}
\Term{skeleton}{骨架}\Term{functor}{函子}与\Term{coskeleton}{余骨架}\Term{functor}{函子}在\Term{semisimplicial-set}{半单纯集}理论中经常使用 [3]。
\SourcePage{44}
\begin{SGAresult}{命题}{7.1.2}
1）\Term{adjunction}{伴随}\Term{morphism}{态射} $\operatorname{sk}_n\to\operatorname{id}$ 是\Term{monomorphism}{单态射}。当 $n\leq m$（相应地 $n\geq m$）时，\Term{canonical}{典范}\Term{morphism}{态射} $\operatorname{sk}_n\operatorname{sk}_m\to\operatorname{sk}_n$（相应地 $\operatorname{cosk}_n\to\operatorname{cosk}_n\operatorname{cosk}_m$）是\Term{isomorphism}{同构}。

2）设 $u:E\to E'$ 是\Term{topos}{拓扑斯}\Term{morphism}{态射}。把\Term{functor}{函子} $u^*$ 延拓到\Term{semisimplicial-object}{半单纯对象}及截断\Term{semisimplicial-object}{半单纯对象}后，它与\Term{functor}{函子} $i_{n!},i_n^*,\operatorname{sk}_n,\operatorname{cosk}_n$ 交换。
\end{SGAresult}
\begin{proof}第一个断言立即由定义得到。证明 2）时，只需注意，由第 I 讲 5.1，\Term{functor}{函子} $i_{n!},\ldots,\operatorname{cosk}_n$ 可用\Term{inductive-limit}{归纳极限}和\Term{finite}{有限}\Term{projective-limit}{投射极限}计算。\end{proof}
\textbf{7.2。一个\Term{homologically-acyclic}{无同调}性引理。}

\textbf{7.2.0。}设 $M$ 是 $E$ 的\Term{abelian-group}{阿贝尔群}。$E$ 的\Term{semisimplicial-object}{半单纯对象} $K$ 以 $M$ 为系数的\Term{homology}{同调}按通常方式定义：先取由 $K$ 生成的自由\Term{semisimple}{半单}纯\Term{abelian-group}{阿贝尔群} $A(K)$，再构造\Term{tensor-product}{张量积} $M\otimes_{\mathbb Z}A(K)$，得到 $E$ 的\Term{semisimple}{半单}纯\Term{abelian-group}{阿贝尔群}。然后考虑相应的\Term{abelian-group}{阿贝尔群}\Term{complex}{复形}（取\Term{boundary-operator}{边界算子}的交错和），并取其\Term{homology}{同调}。\Term{homology-object}{同调对象}记作 $H_i(K,M)$，它们是 $E$ 的\Term{abelian-group}{阿贝尔群}。当 $M$ 是群 $\mathbb Z_E$（由 $E$ 的\Term{final-object}{终对象}生成的自由\Term{abelian-group}{阿贝尔群}）时，简记作 $H_i(K)$。群 $H_i(K)$ 称为 $K$ 的\Term{homology-group}{同调群}。如果对 $E$ 的每个\Term{abelian-group}{阿贝尔群} $M$，群 $H_i(K,M)$（$i>0$）均为零，或者\Term{equivalence}{等价}地，群 $H_i(K,M)$（$i>0$）均为零，就称\Term{semisimplicial-object}{半单纯对象} $K$ 是\Term{homologically-acyclic}{无同调的}。

\Term{homology}{同调}的形成与\Term{topos}{拓扑斯}\Term{morphism}{态射}的\Term{inverse-image}{逆像}交换。

本小节的目的是证明下述引理。
\begin{SGAresult}{引理}{7.2.1}
设 $E$ 是\Term{topos}{拓扑斯}，$K_\bullet,L_\bullet$ 是 $E$ 的两个\Term{semisimplicial-object}{半单纯对象}，$v_\bullet:K_\bullet\to L_\bullet$ 是\Term{semisimplicial-object}{半单纯对象}的\Term{morphism}{态射}，$n$ 是 $\geq0$ 的整数。假设\Term{morphism}{态射} $v$ 具有下列性质：

1）对每个整数 $p<n$（$p\geq0$），\Term{morphism}{态射} $v_p:K_p\to L_p$ 是\Term{isomorphism}{同构}；

2）\Term{morphism}{态射} $v_n:K_n\to L_n$ 是\Term{epimorphism}{满态射}；

3）\Term{canonical}{典范}\Term{morphism}{态射} $K_\bullet\to\operatorname{cosk}_n(K_\bullet)$ 与 $L_\bullet\to\operatorname{cosk}_n(L_\bullet)$ 是\Term{isomorphism}{同构}。

则对每个整数 $p$，\Term{morphism}{态射} $v_p$ 是\Term{epimorphism}{满态射}，而\Term{morphism}{态射} $v_\bullet$ 在\Term{homology-object}{同调对象}上诱导\Term{isomorphism}{同构}。
\end{SGAresult}
\textbf{注记。}证明实际上表明：在\Term{topos}{拓扑斯}中采用适当的\Term{homotopy}{同伦}定义后，\Term{morphism}{态射} $v_\bullet$ 是\Term{homotopy-equivalence}{同伦等价}。
\begin{proof}
可以假设 $E$ 是一个小\Term{site}{景} $C$ 上的\Term{sheaf}{层}\Term{topos}{拓扑斯}（第 IV 讲 1），从而 $E$ 是某个有足够多\Term{fiber}{纤维}\Term{functor}{函子}的\Term{topos}{拓扑斯} $E'$ 的子\Term{topos}{拓扑斯}
\SourcePage{45}
（例如\Term{topos}{拓扑斯} $\widehat C$）。以 $a^*:E'\to E$ 表示\Term{embedding}{嵌入}\Term{morphism}{态射} $E\hookrightarrow E'$ 的\Term{inverse-image}{逆像}\Term{functor}{函子}。设
\[v_\bullet[n]:K_\bullet[n]\longrightarrow L_\bullet[n]\]
是将 $v_\bullet$ 截断到 $n$ 阶而得到的截断\Term{semisimplicial-object}{半单纯对象}的\Term{morphism}{态射}。以 $L_\bullet[n]'$ 表示 $K_\bullet[N]$ 在 $v_\bullet[n]$ 下的像（在 $E'$ 中），以 $v_\bullet[N]':K_\bullet[n]\to L_\bullet[n]'$ 表示 $v_\bullet[n]$ 诱导的\Term{morphism}{态射}。置 $L_!=i_{n*}L_\bullet[n]'$，$K_!=i_{n*}K_\bullet[n]$，$v_!=i_{n*}v_\bullet[n]'$（这里\Term{functor}{函子} $i_{n*}$ 是相对于 $E'$ 而言）。根据 1），\Term{morphism}{态射} $v_!:K_!\to L_!$ 具有引理的性质 1）、2）、3）。此外，根据 2）、3）和 7.1.2，\Term{morphism}{态射} $a^*v_!$ 正是 $v_\bullet$。所以只需对 $v_!$ 证明引理，因而可以假设 $E$ 有足够多\Term{fiber}{纤维}\Term{functor}{函子}；再利用这些\Term{fiber}{纤维}\Term{functor}{函子}，就可归约到 $E$ 为\Term{set}{集合}\Term{category}{范畴}的情形。

首先可以看出，引理关于\Term{morphism}{态射} $v_\bullet$ 的假设在沿任意 $M_\bullet\to L_\bullet$ 作\Term{base-change}{基变换}后仍成立，其中 $M_\bullet$ 是 $n$ 阶\Term{coskeleton}{余骨架}。因此，$v_\bullet$ 在 $L_\bullet$ 的任意基点处的\Term{fiber}{纤维} $P_\bullet$ 形如 $i_{n*}P_\bullet[n]$，其中 $P_\bullet[n]$ 是形如下式的\Term{nonempty}{非空} $n$ 阶截断\Term{complex}{复形}：
\begin{equation}P_n\longrightarrow e\longrightarrow e\longrightarrow e\longrightarrow\cdots\longrightarrow e.\tag{7.2.2}\end{equation}
对每个整数 $i$，以 $\delta^i$ 表示单形 $\Delta^i$ 的\Term{boundary}{边界}\Term{semisimplicial-set}{半单纯集}。每个从 $\delta^i$ 到 $P_\bullet$ 的\Term{morphism}{态射}都可延拓为从 $\Delta^i$ 到 $P_\bullet$ 的\Term{morphism}{态射}。事实上，若 $i\geq n$，由于 $P_\bullet$ 是 $n$ 阶\Term{coskeleton}{余骨架}，这一性质显然成立；若 $i<n$，则由描述 (7.2.2) 也显然成立。由于 $P_\bullet$ 是 Kan \Term{complex}{复形}，$P_\bullet$ 是\Term{homotopy}{同伦}平凡的 [6]。注意\Term{morphism}{态射} $v_\bullet$ 是 Kan 意义下的\Term{kan-fibration}{纤维化} [6]。此\Term{kan-fibration}{纤维化}的\Term{fiber}{纤维}\Term{homotopy}{同伦}平凡，故 $v_\bullet$ 在 $K_\bullet$ 的所有基点处诱导 $K_\bullet$ 与 $L_\bullet$ 的\Term{homotopy-group}{同伦群}的\Term{isomorphism}{同构}。由 Hurewicz 定理，这意味着 $v_\bullet$ 在\Term{homology}{同调}上诱导\Term{isomorphism}{同构} [6]。引理的第二个断言因而得证。为证明第一个断言，注意每个从 $n$ 阶截断\Term{complex}{复形}到 $L_\bullet[n]$ 的\Term{morphism}{态射}，都\Term{lifting}{提升}为到 $K_\bullet[n]$ 的\Term{morphism}{态射}。因而由 3），每个从\Term{complex}{复形}到 $L_\bullet$ 的\Term{morphism}{态射}都\Term{lifting}{提升}为到 $L_\bullet$ 的\Term{morphism}{态射}。所以 $v_\bullet$ 是\Term{surjectivity}{满射}。
\end{proof}
\textbf{7.3。\Term{hypercovering}{超覆盖}。}

\textbf{7.3.0。}设 $C$ 是属于\Term{universe}{宇宙} $\mathcal U$ 的\Term{site}{景}，其中\Term{fiber-product}{纤维积}及两个\Term{object}{对象}的乘积均\Term{representable-functor}{可表示}。以 $\widehat C$ 表示 $C$ 上的 $\mathcal U$-\Term{presheaf}{预层}\Term{topos}{拓扑斯}，以 $C^{\sim}$ 表示 $C$ 上的 $\mathcal U$-\Term{sheaf}{层}\Term{topos}{拓扑斯}。如果 $\widehat C$ 的\Term{object}{对象} $K$ \Term{isomorphism}{同构}于一些\Term{representable-functor}{可表示}\Term{presheaf}{预层}的\Term{direct-sum}{直和}，就称它\Term{semi-representable}{半可表示}。

设 $\operatorname{SR}(C)$ 是 $\widehat C$ 中由\Term{semi-representable}{半可表示}\Term{object}{对象}定义的\Term{full-subcategory}{满子范畴}。在\Term{category}{范畴} $\operatorname{SR}(C)$ 中，以\Term{nonempty}{非空}\Term{finite-category}{有限范畴}为指标的\Term{projective-limit}{投射极限}均\Term{representable-functor}{可表示}，且包含\Term{functor}{函子} $\operatorname{SR}(C)\hookrightarrow\widehat C$ 与这些\Term{finite}{有限}\Term{projective-limit}{投射极限}交换。由已作出的一个注记可知：若 $K_\bullet$ 是
\SourcePage{46}
$\widehat C$ 中截断到 $n$ 阶的\Term{semisimplicial-object}{半单纯对象}，且其中所有\Term{object}{对象}均\Term{semi-representable}{半可表示}，则延拓 $i_n^+(K_\bullet)$ 也具有此性质。因此，若 $K_\bullet$ 是 $\widehat C$ 的\Term{semisimplicial-object}{半单纯对象}，其中所有\Term{object}{对象}均\Term{semi-representable}{半可表示}，则对每个 $n$，\Term{object}{对象} $\operatorname{cosk}_n(K_\bullet)$ 也具有此性质。

\textbf{7.3.1。定义与记号。}

\textbf{7.3.1.1。}设 $p>0$ 是整数。若 $\widehat C$ 的\Term{semisimplicial-object}{半单纯对象} $K_\bullet$ 具有下列性质，就称它为 $C$ 的\Term{hypercovering}{型为 $p$ 的超覆盖}，或在不致混淆时简称型为 $p$ 的\Term{hypercovering}{超覆盖}：

HR1）对每个整数 $n\geq0$，$K_n$ \Term{semi-representable}{半可表示}。

HR2）\Term{canonical}{典范}\Term{morphism}{态射} $K_\bullet\to\operatorname{cosk}_p(K_\bullet)$ 是\Term{isomorphism}{同构}。

HR3）对每个整数 $n\geq0$，\Term{presheaf}{预层}的\Term{canonical}{典范}\Term{morphism}{态射} $K_{n+1}\to(\operatorname{cosk}_n(K_\bullet))_{n+1}$ 是\Term{presheaf}{预层}的\Term{covering}{覆盖}\Term{isomorphism}{同构}（第 II 讲 6.2）。\Term{presheaf}{预层}的\Term{canonical}{典范}\Term{morphism}{态射} $K_0\to e$ 也是\Term{presheaf}{预层}的\Term{covering}{覆盖}\Term{morphism}{态射}，其中 $e$ 是 $\widehat C$ 的\Term{final-object}{终对象}。

\textbf{7.3.1.2。}对每个 $q\geq p$，型为 $p$ 的\Term{hypercovering}{超覆盖}也是型为 $q$ 的\Term{hypercovering}{超覆盖}。若\Term{semisimplicial-object}{半单纯对象} $K_\bullet$ 具有性质 HR1）与 HR3），便称它为（$C$ 的）\Term{hypercovering}{超覆盖}。

\textbf{7.3.1.3。}以下\Term{category}{范畴}记作 $\operatorname{HR}_p$（相应地 $\operatorname{HR}$），称为\Term{hypercovering}{型为 $p$ 的超覆盖范畴}（相应地，\Term{hypercovering}{超覆盖}\Term{category}{范畴}）：

a）$\operatorname{HR}_p$（相应地 $\operatorname{HR}$）的\Term{object}{对象}是型为 $p$ 的\Term{hypercovering}{超覆盖}（相应地，\Term{hypercovering}{超覆盖}）。

b）若 $K_\bullet,L_\bullet$ 是 $\operatorname{HR}_p$（相应地 $\operatorname{HR}$）的两个\Term{object}{对象}，则源为 $K_\bullet$、靶为 $L_\bullet$ 的\Term{morphism}{态射}，就是 $\widehat C$ 中\Term{semisimplicial-object}{半单纯对象}的\Term{morphism}{态射} $v_\bullet:K_\bullet\to L_\bullet$。

\textbf{7.3.1.4。}设 $C$ 是\Term{fiber-product}{纤维积}均\Term{representable-functor}{可表示}的\Term{site}{景}，$X$ 是 $C$ 的\Term{object}{对象}。$X$ 的\Term{hypercovering}{超覆盖}（相应地，型为 $p$ 的\Term{hypercovering}{超覆盖}），指\Term{site}{景} $C/X$ 的\Term{hypercovering}{超覆盖}（相应地，型为 $p$ 的\Term{hypercovering}{超覆盖}）。同样定义 $X$ 的\Term{hypercovering}{超覆盖}\Term{category}{范畴}（相应地，型为 $p$ 的\Term{hypercovering}{超覆盖}\Term{category}{范畴}）。有一列包含\Term{functor}{函子}
\[\cdots\operatorname{HR}_p\hookrightarrow\operatorname{HR}_{p+1}\hookrightarrow\cdots\hookrightarrow\operatorname{HR}.\]
这些\Term{functor}{函子}都是\Term{fully-faithful}{全忠实}的。以 $\operatorname{HR}_\infty$ 表示\Term{category}{范畴} $\operatorname{HR}_p$ 的\Term{inductive-limit}{归纳极限}。

\textbf{7.3.1.6。}对每个整数 $n\geq0$，以 $\Delta_n^c$ 表示\Term{set}{集合}\Term{category}{范畴}中\Term{dimension}{维数}为 $n$ 的标准\Term{semisimplicial-object}{半单纯对象}，即 $\boldsymbol\Delta$ 上由\Term{ordered-set}{有序集} $[0,n]$ 表示的\Term{functor}{函子}。以 $\Delta_n^c$ 表示 $C$ 上取值于\Term{semisimplicial-set}{半单纯集}\Term{category}{范畴}、以 $\Delta_n$ 为值的常\Term{presheaf}{预层}。因此\Term{presheaf}{预层} $\Delta_n^c$ 是\Term{semisimplicial-set}{半单纯集}合\Term{presheaf}{预层}。从 $\Delta_0$ 到 $\Delta_1$ 的两个\Term{canonical}{典范}\Term{embedding}{嵌入}
\SourcePage{47}
定义\Term{semisimplicial-presheaf}{半单纯预层}的两个\Term{morphism}{态射}
\[\begin{tikzcd}\Delta_0^c\arrow[r,shift left,"e_0"]\arrow[r,shift right,"e_1"']&\Delta_1^c,\end{tikzcd}\]
从而对每个\Term{semisimplicial-presheaf}{半单纯预层} $K_\bullet$ 定义两个\Term{canonical}{典范}\Term{embedding}{嵌入}
\[\begin{tikzcd}K_\bullet\arrow[r,shift left,"e_0"]\arrow[r,shift right,"e_1"']&K_\bullet\times\Delta_1^c.\end{tikzcd}\]
\Term{semisimplicial-presheaf}{半单纯预层}的两个\Term{morphism}{态射} $u_0,u_1:K_\bullet\rightrightarrows L_\bullet$ 称为\Term{homotopic-morphisms}{同伦的态射}，若存在\Term{morphism}{态射} $v:K_\bullet\times\Delta_1^c\to L_\bullet$，使图
\[\begin{tikzcd}K_\bullet\arrow[r,"e_i"]\arrow[dr,"u_i"']&K_\bullet\times\Delta_1^c\arrow[d,"v"]\\&L_\bullet\end{tikzcd}\qquad i=0,1\]
交换。\Term{morphism}{态射} $v$ 称为连接 $u_0$ 与 $u_1$ 的\Term{homotopy}{同伦}。“$u_0$ 与 $u_1$ 是从 $K_\bullet$ 到 $L_\bullet$ 的两个\Term{homotopic-morphisms}{同伦的态射}”这一关系一般不是\Term{equivalence-relation}{等价关系}。不过，由此关系生成的\Term{equivalence-relation}{等价关系}与\Term{morphism}{态射}的复合相容。因此，将\Term{morphism}{态射}按\Term{homotopy}{同伦}关系生成的\Term{equivalence-relation}{等价关系}取商，就可以定义模\Term{homotopy}{同伦}的\Term{semisimplicial-presheaf}{半单纯预层}\Term{category}{范畴}。

\textbf{7.3.1.7。}这样定义\Term{category}{范畴} $\mathscr{HR}_p,\mathscr{HR}_\infty,\mathscr{HR}$，即\Term{hypercovering}{模同伦的超覆盖范畴}。
\begin{SGAresult}{定理}{7.3.2}
设\Term{site}{景} $C$ 满足 7.3.0 的条件。

1）\Term{category}{范畴} $\mathscr{HR}_p^{\circ}$（相应地 $\mathscr{HR}^{\circ}$）是\Term{filtered-category}{滤过}的（第 I 讲 2.7）。

2）设 $K_\bullet$ 是 $\operatorname{HR}_p$（相应地 $\operatorname{HR}$）的\Term{object}{对象}，$n$ 是满足 $0\leq n\leq p$（相应地 $0\leq n$）的整数，$u:X\to K_n$ 是\Term{presheaf}{预层}的\Term{covering}{覆盖}\Term{morphism}{态射}。存在 $\operatorname{HR}_p$（相应地 $\operatorname{HR}$）的\Term{object}{对象} $L_\bullet$ 及\Term{morphism}{态射} $f:L_\bullet\to K_\bullet$，使\Term{morphism}{态射} $f_n:L_n\to K_n$ 经由 $u$ \Term{resolution}{分解}。
\SourcePage{48}
3）\Term{hypercovering}{超覆盖}的\Term{adjunction}{伴随}\Term{semisimplicial-sheaf}{半单纯层}（第 II 讲 3.5）在严格正次数上是\Term{homologically-acyclic}{无同调}的（7.2.0）。第零\Term{homology}{同调}\Term{sheaf}{层}\Term{isomorphism}{同构}于常\Term{presheaf}{预层} $\mathbb Z$ 的\Term{associated-sheaf}{伴随层}。
\end{SGAresult}
\begin{proof}[3）的证明]
设 $K_\bullet^{\sim}$ 是 $K_\bullet$ 的\Term{adjunction}{伴随}\Term{semisimplicial-sheaf}{半单纯层}。“\Term{associated-sheaf}{伴随层}”\Term{functor}{函子}与\Term{functor}{函子} $\operatorname{cosk}_n$ 交换（7.1.2，2））。所以\Term{semisimplicial-sheaf}{半单纯层} $K_\bullet^{\sim}$ 具有如下性质：

a）\Term{canonical}{典范}\Term{morphism}{态射} $K_0\to e$ 是\Term{sheaf}{层}的\Term{epimorphism}{满态射}。

b）对每个整数 $n\geq0$，\Term{canonical}{典范}\Term{morphism}{态射} $K_{n+1}\to(\operatorname{cosk}_n(K_\bullet))_{n+1}$ 是\Term{sheaf}{层}的\Term{epimorphism}{满态射}。

对 $n\geq0$，置 $L_{\bullet n}=\operatorname{cosk}_n K_\bullet^{\sim}$。于是对每个 $n$，有一列\Term{semisimplicial-sheaf}{半单纯层}的\Term{morphism}{态射}
\[K_\bullet^{\sim}\xrightarrow{u_n}L_{\bullet n}\xrightarrow{v_n}L_{\bullet,n-1}\longrightarrow\cdots\longrightarrow L_{\bullet0}\xrightarrow{v_0}e_\bullet,\]
且对每个 $0\leq j\leq n$，$v_j$ 具有 7.2.1 的性质。此外，\Term{morphism}{态射} $u_n$ 在次数 $\leq n$ 的\Term{set-component}{分量}上诱导\Term{isomorphism}{同构}。对 $n$ 归纳，由 2.1 就得到 $K_\bullet$ 是\Term{homologically-acyclic}{无同调}的。

为证明 1）与 2），先引入一个术语。
\end{proof}
\begin{SGAresult}{定义}{7.3.3}
设 $f:X_\bullet\to Y_\bullet$ 是\Term{semisimplicial-presheaf}{半单纯预层}的\Term{morphism}{态射}。若满足下列条件，就称 $f$ 是\Term{special-morphism}{型为 $p$ 的特殊态射}（相应地，\Term{special-morphism}{特殊态射}）：

1）\Term{object}{对象} $X_\bullet,Y_\bullet$ 与其 $p$ 阶\Term{coskeleton}{余骨架}\Term{canonical-isomorphism}{典范同构}，且 $\operatorname{cosk}_p(f)=f$（相应地，不对 $f,X_\bullet,Y_\bullet$ 加条件）。

2）对每个满足 $0\leq n\leq p$（相应地 $0\leq n$）的整数 $n$，下图中的\Term{morphism}{态射} $\Phi_{n+1}$ 是\Term{covering}{覆盖}\Term{morphism}{态射}：
\[\begin{tikzcd}[column sep=large,row sep=large]
X_{n+1}\arrow[rr,"f_{n+1}"]\arrow[dd]\arrow[dr,"\Phi_{n+1}"]&&Y_{n+1}\arrow[dd]\\
&P_{n+1}\arrow[ur]\arrow[dl]&\\
(\operatorname{cosk}_nX_\bullet)_{n+1}\arrow[rr,"(\operatorname{cosk}_nf)_{n+1}"']&&(\operatorname{cosk}_nY_\bullet)_{n+1}.
\end{tikzcd}\]
（竖\Term{arrow}{箭头}由\Term{canonical}{典范}\Term{morphism}{态射} $X_\bullet\to\operatorname{cosk}_nX_\bullet$ 及 $Y_\bullet\to\operatorname{cosk}_nY_\bullet$ 定义。\Term{object}{对象} $P_{n+1}$ 是\Term{fiber-product}{纤维积}，$\Phi_{n+1}$ 是使此图交换的唯一\Term{arrow}{箭头}。）

3）\Term{morphism}{态射} $f_0:X_0\to Y_0$ 是\Term{covering}{覆盖}\Term{morphism}{态射}。
\end{SGAresult}
\SourcePage{49}
若\Term{canonical}{典范}\Term{morphism}{态射} $K_\bullet\to e_\bullet$（$e_\bullet$ 是终\Term{semisimplicial-object}{半单纯对象}）是型为 $p$ 的\Term{special-morphism}{特殊态射}（相应地，\Term{special-morphism}{特殊态射}），即 $K_\bullet$ 满足 7.3.1.1 的条件 HR2）、HR3）（相应地 HR3）），就称\Term{semisimplicial-presheaf}{半单纯预层} $K_\bullet$ 是型为 $p$ 的特殊\Term{object}{对象}（相应地，特殊\Term{object}{对象}）。
\begin{SGAresult}{引理}{7.3.4}
1）两个\Term{special-morphism}{特殊态射}（相应地，型为 $p$ 的\Term{special-morphism}{特殊态射}）的复合仍是\Term{special-morphism}{特殊态射}（相应地，型为 $p$ 的\Term{special-morphism}{特殊态射}）。

2）设 $K_\bullet$ 是特殊\Term{semisimplicial-presheaf}{半单纯预层}（相应地，型为 $p$ 的特殊\Term{semisimplicial-presheaf}{半单纯预层}），$f:X_\bullet\to Y_\bullet$ 是\Term{special-morphism}{特殊态射}（相应地，型为 $p$ 的\Term{special-morphism}{特殊态射}），$u:K_\bullet\to Y_\bullet$ 是\Term{complex}{复形}的\Term{morphism}{态射}。则\Term{fiber-product}{纤维积} $P_\bullet=K_\bullet\times_{Y_\bullet}X_\bullet$ 是特殊\Term{object}{对象}（相应地，型为 $p$ 的特殊\Term{object}{对象}）。
\end{SGAresult}
此引理的证明留给读者。

\textbf{7.3.5。7.3.2 中断言 2）的证明。}首先可假设 $X$ \Term{semi-representable}{半可表示}。对每个\Term{semisimplicial-object}{半单纯对象} $L_\bullet$，以 $\operatorname{Hom}_{(X)}(L_\bullet,K_\bullet)$ 表示附带\Term{resolution}{分解} $L_n\to X\xrightarrow{u}K_n$ 的\Term{semisimplicial-object}{半单纯对象}\Term{morphism}{态射}所成的\Term{set}{集合}。\Term{functor}{函子} $\operatorname{Hom}_{(X)}(\cdot,K_\bullet)$ \Term{representable-functor}{可表示}。事实上，它把\Term{inductive-limit}{归纳极限}变为\Term{projective-limit}{投射极限}，且\Term{category}{范畴} $\boldsymbol\Delta(\widehat C)$ 是\Term{topos}{拓扑斯}。以 $P_\bullet$ 表示代表此\Term{functor}{函子}的\Term{object}{对象}。对 $\widehat C$ 的\Term{object}{对象} $Z$，同样以 $j_n^+(Z)$ 表示\Term{category}{范畴} $\boldsymbol\Delta\widehat C$ 上\Term{functor}{函子} $L_\bullet\mapsto\operatorname{Hom}(L_n,Z)$ 的代表\Term{object}{对象}；利用\Term{projective-limit}{投射极限}所作的此\Term{object}{对象}的构造，已在第 III 讲 1.1 中明确给出。这里可见，此构造只涉及\Term{finite}{有限}\Term{projective-limit}{投射极限}，因而立即知道 $j_n^+$ 将\Term{semi-representable}{半可表示}\Term{presheaf}{预层}变为各\Term{set-component}{分量}均\Term{semi-representable}{半可表示}的\Term{semisimplicial-presheaf}{半单纯预层}，并将\Term{covering}{覆盖}\Term{morphism}{态射} $Z'\to Z$ 变为\Term{special-morphism}{特殊态射}（7.3.3）。根据 $P_\bullet$ 的定义，有\Term{cartesian}{笛卡尔}方图
\[\begin{tikzcd}P_\bullet\arrow[r]\arrow[d]&j_n^+(X)\arrow[d,"j_n^+(u)"]\\K_\bullet\arrow[r]&j_n^+(K_n).
\end{tikzcd}\]
由刚才所说，\Term{semisimplicial-object}{半单纯对象} $j_n^+(X)$ 与 $j_n^+(K_n)$ 均\Term{semi-representable}{半可表示}，而\Term{morphism}{态射} $j_n^+(u)$ 是型为 $n$ 的\Term{special-morphism}{特殊态射}。于是只需应用 3.4 及 7.3.0 的一系列事实，即得结论。

\textbf{7.3.6.0。}设 $M_\bullet,N_\bullet$ 是两个\Term{semisimplicial-presheaf}{半单纯预层}。\Term{functor}{函子}
\[L_\bullet\longmapsto\operatorname{Hom}(L_\bullet\times M_\bullet,N_\bullet)\]
\Term{representable-functor}{可表示}。其代表\Term{semisimplicial-presheaf}{半单纯预层}记作 $\mathscr Hom_\bullet(M_\bullet,N_\bullet)$。
\SourcePage{50}
此\Term{object}{对象}的 $n$ 次\Term{set-component}{分量}\Term{presheaf}{预层}记作 $\mathscr Hom_n(M_\bullet,N_\bullet)$。有\Term{canonical-isomorphism}{典范同构}
\[\mathscr Hom_n(M_\bullet,N_\bullet)\xrightarrow{\sim}\mathscr Hom_0(M_\bullet\times\Delta_n^c,N_\bullet).\]
\begin{SGAresult}{引理}{7.3.6}
设
\[\begin{tikzcd}M_\bullet\arrow[r]&N_\bullet\\
\operatorname{sk}_n\Delta_{n+1}^c\arrow[u,hook]\arrow[r,hook,"u"']&\Delta_{n+1}^c\arrow[u]
\end{tikzcd}\qquad\text{（$u$ 是典范嵌入）}\]
是\Term{semisimplicial-presheaf}{半单纯预层}的\Term{cocartesian}{余笛卡尔}图。则对每个\Term{hypercovering}{超覆盖} $L_\bullet$，\Term{morphism}{态射}
\[\mathscr Hom_0(N_\bullet,L_\bullet)\longrightarrow\mathscr Hom_0(M_\bullet,L_\bullet)\]
是\Term{covering}{覆盖}\Term{morphism}{态射}。
\end{SGAresult}
\begin{proof}有\Term{cartesian}{笛卡尔}图
\[\begin{tikzcd}[column sep=large]
\mathscr Hom_0(N_\bullet,L_\bullet)\arrow[r]\arrow[d]&\mathscr Hom_0(M_\bullet,L_\bullet)\arrow[d]\\
\mathscr Hom_0(\Delta_{n+1}^c,L_\bullet)\arrow[r,"(*)"']&\mathscr Hom_0(\operatorname{sk}_n\Delta_{n+1}^c,L_\bullet).
\end{tikzcd}\]
所以只需证明\Term{morphism}{态射} $(*)$ 是\Term{covering}{覆盖}\Term{morphism}{态射}。而它\Term{isomorphism}{同构}于\Term{morphism}{态射}
\[L_{n+1}\longrightarrow(\operatorname{cosk}_nL_\bullet)_{n+1},\]
后者由假设是\Term{covering}{覆盖}\Term{morphism}{态射}，证毕。\end{proof}
\textbf{7.3.7。7.3.2 中断言 1）的证明。}$\operatorname{HR}_p$（相应地 $\operatorname{HR}$）中两个\Term{object}{对象}的乘积仍属于 $\operatorname{HR}_p$（相应地 $\operatorname{HR}$）（7.3.4）。为证明\Term{category}{范畴} $\mathscr{HR}_p^\circ$（相应地 $\mathscr{HR}^\circ$）\Term{filtered}{滤过}，只需证明：给定 $\operatorname{HR}_p$（相应地 $\operatorname{HR}$）的两个\Term{morphism}{态射} $u_0,u_1:K_\bullet\rightrightarrows L_\bullet$，存在 $\operatorname{HR}_p$（相应地 $\operatorname{HR}$）中的\Term{morphism}{态射} $v:M_\bullet\to K_\bullet$，使\Term{morphism}{态射} $u_0v,u_1v:M_\bullet\rightrightarrows L_\bullet$
\SourcePage{51}
\Term{homotopy}{同伦}，即存在 $w:M_\bullet\times\Delta_1^c\to L_\bullet$，使下图交换：
\[\begin{tikzcd}M_\bullet\arrow[r,"e_i"]\arrow[d,"v"']&M_\bullet\times\Delta_1^c\arrow[d,"w"]\\
K_\bullet\arrow[r,"u_i"']&L_\bullet\end{tikzcd}\qquad i=0,1.\tag{x}\]
对 $\widehat C$ 的每个\Term{semisimplicial-object}{半单纯对象} $M_\bullet$（不要求为\Term{hypercovering}{超覆盖}），设 $F(M_\bullet)$ 是满足图 (x) 交换的二元组 $(v,w)$ 的\Term{set}{集合}，其中 $v:M_\bullet\to K_\bullet$，$w:M_\bullet\times\Delta_1^c\to L_\bullet$。\Term{functor}{函子} $M_\bullet\mapsto F(M_\bullet)$ 关于 $M_\bullet$ \Term{contravariant}{反变}，并把\Term{inductive-limit}{归纳极限}变为\Term{projective-limit}{投射极限}，因此由某个\Term{semisimplicial-object}{半单纯对象} $F_\bullet$ 表示。\Term{object}{对象} $F_\bullet$ 显然位于\Term{cartesian}{笛卡尔}图
\[\begin{tikzcd}F_\bullet\arrow[r]\arrow[d]&\mathscr Hom_\bullet(\Delta_1^c,L_\bullet)\arrow[d,"\pi"]\\
K_\bullet\arrow[r,"{u_0,u_1}"']&L_\bullet\times L_\bullet,
\end{tikzcd}\]
其中 $\pi$ 由 $\Delta_0^c$ 到 $\Delta_1^c$ 的两个包含映射定义。因此，为证明所述断言，只需证明 $F_\bullet$ 是 $\operatorname{HR}_p$（相应地 $\operatorname{HR}$）的\Term{object}{对象}；为此，由 7.3.4 的 2）及 7.3.0 的一系列事实，只需证明：

a）$\mathscr Hom_\bullet(\Delta_1^c,L_\bullet)$ \Term{semi-representable}{半可表示}；

b）\Term{morphism}{态射} $\pi$ 是型为 $p$ 的\Term{special-morphism}{特殊态射}（相应地，\Term{special-morphism}{特殊态射}）。

首先立即验证 $\mathscr Hom_\bullet(\Delta_1^c,L_\bullet)$ 是\Term{semi-representable}{半可表示}的\Term{semisimplicial-object}{半单纯对象}。事实上，其各\Term{set-component}{分量}通过对 $L_n$ 作\Term{finite}{有限}\Term{projective-limit}{投射极限}计算，所以\Term{semi-representable}{半可表示}。现在验证 b）。首先可证明：当 $L_\bullet$ 是型为 $p$ 的特殊\Term{object}{对象}时，\Term{morphism}{态射}
\[\mathscr Hom_\bullet(\Delta_1^c,L_\bullet)\longrightarrow\operatorname{cosk}_p\mathscr Hom_\bullet(\Delta_1^c,L_\bullet)\]
是\Term{isomorphism}{同构}；于是可验证 7.3.3 的性质 1）。其次，\Term{morphism}{态射}
\[\pi_0:\mathscr Hom_0(\Delta_1^c,L_\bullet)\longrightarrow L_0\times L_0\]
\Term{isomorphism}{同构}于\Term{canonical}{典范}\Term{morphism}{态射} $L_1\xrightarrow{(d_0,d_1)}L_0\times L_0$，后者由假设是\Term{covering}{覆盖}\Term{morphism}{态射}（$L_\bullet$ 是\Term{hypercovering}{超覆盖}）。
\SourcePage{52}
所以尚须验证 7.3.3 的性质 2），即对每个 $n$，在图
\[\begin{tikzcd}[column sep=large,row sep=large]
\mathscr Hom_{n+1}(\Delta_1^c,L_\bullet)\arrow[rr,"\pi_{n+1}"]\arrow[dd]\arrow[dr,"\Phi_{n+1}"]&&L_{n+1}\times L_{n+1}\arrow[dd]\\
&P_{n+1}\arrow[ur]\arrow[dl]&\\
(\operatorname{cosk}_n\mathscr Hom_\bullet(\Delta_1^c,L_\bullet))_{n+1}\arrow[rr,"(\operatorname{cosk}_n\pi)_{n+1}"']&&(\operatorname{cosk}_nL_\bullet)_{n+1}\times(\operatorname{cosk}_nL_\bullet)_{n+1}
\end{tikzcd}\]
（$P_{n+1}$ 是\Term{fiber-product}{纤维积}）中，\Term{morphism}{态射} $\Phi_{n+1}$ 是\Term{covering}{覆盖}\Term{morphism}{态射}。而简单地追踪\Term{adjunction}{伴随}\Term{functor}{函子}便得到
\begin{align*}
P_{n+1}&\xrightarrow{\sim}\mathscr Hom_0(\operatorname{sk}_{n+1}(\Delta_{n+1}^c\times\Delta_1^c),L_\bullet),\\
\mathscr Hom_{n+1}(\Delta_1^c,L_\bullet)&\xrightarrow{\sim}\mathscr Hom_0(\Delta_{n+1}^c\times\Delta_1^c,L_\bullet),
\end{align*}
并知\Term{morphism}{态射} $\Phi_{n+1}$ \Term{isomorphism}{同构}于由\Term{embedding}{嵌入}
\[\operatorname{sk}_{n+1}(\Delta_{n+1}^c\times\Delta_1^c)\hookrightarrow\Delta_{n+1}^c\times\Delta_1^c\]
产生的\Term{morphism}{态射}
\[\mathscr Hom_0(\Delta_{n+1}^c\times\Delta_1^c,L_\bullet)\longrightarrow\mathscr Hom_0(\operatorname{sk}_{n+1}(\Delta_{n+1}^c\times\Delta_1^c),L_\bullet).\]
而在 $\Delta_{n+1}^c\times\Delta_1^c$ 中存在\Term{subobject}{子对象}列
\[\operatorname{sk}_{n+1}(\Delta_{n+1}^c\times\Delta_1^c)=M_0\hookrightarrow M_1\hookrightarrow\cdots\hookrightarrow M_k=\Delta_{n+1}^c\times\Delta_1^c,\]
使对每个 $0\leq i\leq k$，\Term{morphism}{态射} $M_i\hookrightarrow M_{i+1}$ 都可置于\Term{cocartesian}{余笛卡尔}图
\[\begin{tikzcd}M_i\arrow[r,hook]&M_{i+1}\\
\operatorname{sk}_{n+1}\Delta_{n+2}^c\arrow[u]\arrow[r,hook]&\Delta_{n+2}^c\arrow[u]
\end{tikzcd}\]
中。（逐个添加 $\Delta_{n+1}\times\Delta_1$ 中\Term{dimension}{维数}为 $n+2$ 的非\Term{degeneration}{退化}单形。）由 7.3.6，\Term{morphism}{态射} $\Phi_{n+1}$ 是\Term{covering}{覆盖}\Term{morphism}{态射}的复合，因而自身也是\Term{covering}{覆盖}\Term{morphism}{态射}。7.3.2 的证明至此完成。

\textbf{7.4。\Term{isomorphism}{同构}定理。}
\SourcePage{53}
\textbf{7.4.0。}设 $F$ 是满足条件 7.3.0 的\Term{site}{景} $C$ 上的\Term{abelian-group}{阿贝尔群}\Term{presheaf}{预层}。对每个\Term{hypercovering}{超覆盖} $K_\bullet$，置
\[H^q(K_\bullet,F)=H^q(\operatorname{Hom}_{\widehat C}(K_\bullet,F)).\]
$H^q(K_\bullet,F)$ 在 $C$ 上的阿贝尔\Term{presheaf}{预层}\Term{category}{范畴}上构成 $\delta$-\Term{functor}{函子}，但一般不是\Term{functor}{函子} $H^0(K_\bullet,F)$ 的\Term{derived-functor}{导出函子}。

置
\begin{equation}\overset{r}{\check H}{}^q(C,F)=\varinjlim_{\operatorname{HR}_r}H^q(\cdot,F)\quad\text{（$r$ 可为无穷）},\tag{7.4.0.1}\end{equation}
以及
\begin{equation}\breve H^q(C,F)=\varinjlim_{\operatorname{HR}}H^q(\cdot,F).\tag{7.4.0.2}\end{equation}
$\overset{r}{\check H}{}^q(C,F)$（相应地 $\breve H^q(C,F)$）仍构成 $\delta$-\Term{functor}{函子}，因为\Term{category}{范畴} $\mathscr{HR}_p^\circ$（相应地 $\mathscr{HR}^\circ$）是\Term{filtered-category}{滤过的}（7.3.2）。由于\Term{category}{范畴} $\mathscr{HR}_p$ 未必是 $\mathcal U$-\Term{category}{范畴}，这些 $\delta$-\Term{functor}{函子}取值于某个适当\Term{universe}{宇宙} $\mathcal V$ 中的 $\mathcal V$-\Term{abelian-group}{阿贝尔群}\Term{category}{范畴}。

现在假设\Term{presheaf}{预层} $F$ 是\Term{sheaf}{层}。由于\Term{hypercovering}{超覆盖}的\Term{adjunction}{伴随}\Term{semisimplicial-sheaf}{半单纯层}是\Term{homologically-acyclic}{无同调}的（7.3.2），有关于 $F$ 和 $K_\bullet$ \Term{functor}{函子}性的\Term{spectral-sequence}{谱序列}
\begin{equation}H^p(K_\bullet,\mathscr H^q(F))\Longrightarrow H^{p+q}(C^{\sim},F).\tag{7.4.0.3}\end{equation}
于是取\Term{limit}{极限}，得到\Term{spectral-sequence}{谱序列}
\begin{equation}\begin{aligned}\overset{r}{\check H}{}^p(C,\mathscr H^q(F))&\Longrightarrow H^{p+q}(C^{\sim},F),\\
\breve H^p(C,\mathscr H^q(F))&\Longrightarrow H^{p+q}(C^{\sim},F).
\end{aligned}\tag{7.4.0.4}\end{equation}
\begin{SGAresult}{定理}{7.4.1}
设\Term{site}{景} $C$ 满足条件 7.3.0，$F$ 是 $C$ 上的\Term{abelian-sheaf}{阿贝尔层}。

1）\Term{spectral-sequence}{谱序列} (7.4.0.4) 定义\Term{isomorphism}{同构}
\[\overset{r}{\check H}{}^q(C,F)\xrightarrow{\sim}H^q(C^{\sim},F),\qquad q\leq r+1,\]
以及\Term{monomorphism}{单态射}
\[\overset{r}{\check H}{}^{r+2}(C,F)\longrightarrow H^{r+2}(C^{\sim},F).\]
\SourcePage{54}
2）有 $\delta$-\Term{functor}{函子}的\Term{isomorphism}{同构}
\[\overset{\infty}{\check H}{}^q(C,F)\xrightarrow{\sim}\breve H^q(C,F)\xrightarrow{\sim}H^q(C^{\sim},F)\qquad\text{（对所有 $q$）}.\]
3）设 $G$ 是\Term{abelian-group}{阿贝尔群}\Term{presheaf}{预层}，$F$ 是其\Term{associated-sheaf}{伴随层}。存在\Term{functor}{函子}的\Term{isomorphism}{同构}
\[\overset{r}{\check H}{}^q(C,G)\xrightarrow{\sim}H^q(C^{\sim},F),\qquad q\leq r-1,\]
以及\Term{monomorphism}{单态射}
\[\overset{r}{\check H}{}^r(C,G)\xrightarrow{\sim}H^q(C^{\sim},F).\]
当 $G$ 是\Term{separated}{分离}\Term{presheaf}{预层}时，此最后一个\Term{morphism}{态射}是\Term{isomorphism}{同构}，而且存在\Term{monomorphism}{单态射}
\[\overset{r}{\check H}{}^{r+1}(C,G)\longrightarrow H^{r+1}(C^{\sim},F).\]
4）有\Term{isomorphism}{同构}
\[\overset{\infty}{\check H}{}^q(C,G)\xrightarrow{\sim}\breve H^q(G,G)\xrightarrow{\sim}H^q(C^{\sim},F)\qquad\text{（所有 $q$）}.\]
\end{SGAresult}
\begin{proof}只需证明：若阿贝尔\Term{presheaf}{预层} $N$ 的\Term{associated-sheaf}{伴随层}为零，则当 $q\leq r$ 时，$\overset{r}{\check H}{}^q(C,N)=0$（相应地，对所有 $q$，$\breve H^q(C,N)=0$）；利用 3.2 可立即做到这一点。\end{proof}
\begin{SGAresult}{注记}{7.4.2}
1）注意，对满足 7.3.0 条件的\Term{site}{景}，型为 $0$ 的\Term{hypercovering}{超覆盖}就是通常的\Term{covering}{覆盖}，而\Term{functor}{函子} $\overset{\infty}{\check H}{}^q$ 正是 (2.4.5.1) 中引入的\Term{functor}{函子} $\check H^q$。

2）设 $C$ 是\Term{finite}{有限}\Term{projective-limit}{投射极限}均\Term{representable-functor}{可表示}的\Term{site}{景}，且对 $C$ 的每个\Term{object}{对象} $X$ 及每个\Term{covering}{覆盖}族 $(X_i\to X)$，$i\in I$，都存在 $i_0\in I$，使 $X_{i_0}\to X$ 是\Term{covering}{覆盖}\Term{morphism}{态射}。则可证明：各\Term{set-component}{分量}均\Term{representable-functor}{可表示}的型为 $p$ 的\Term{hypercovering}{超覆盖}在 $\mathscr{HR}_p$ 中\Term{cofinal}{共尾}。同样，满足存在整数 $p$ 使 $K_\bullet$ 型为 $p$、且各\Term{set-component}{分量}\Term{representable-functor}{可表示}的\Term{hypercovering}{超覆盖} $K_\bullet$，在 $\mathscr{HR}$ 中\Term{cofinal}{共尾}。（可以仿照 3.5 证明。）因此在所考虑的情形，用这些\Term{set-component}{分量}\Term{representable-functor}{可表示}的\Term{hypercovering}{超覆盖}就足以计算\Term{presheaf}{预层}的\Term{associated-sheaf}{伴随层}的\Term{cohomology}{上同调}。
\end{SGAresult}
\SourcePage{55}
\SGAsection{8}{附录：局部归纳极限（P. Deligne 撰）}
撰稿人建议读者原则上略过本附录。本附录介绍一种技术，有时能把因存在点而显然成立的断言推广到没有\Term{enough-points}{足够多点}的\Term{topos}{拓扑斯}。这种技术可给出 D. Lazard 定理的\Term{sheaf}{层}版本；该定理断言，一个环上的\Term{flat-module}{平坦模}是\Term{finite-type}{有限型}\Term{free-module}{自由模}的\Term{inductive-limit}{归纳极限}。

\textbf{8.1。\Term{locally-filtered}{局部滤过范畴}。}

\textbf{8.1.0。}若 $\mathscr B$ 是\Term{category}{范畴}，$p:\mathscr A\to\mathscr B$ 是 $\mathscr B$ 上的\Term{category}{范畴}，就使用如下记号：

对 $U\in\operatorname{Ob}\mathscr B$，$\mathscr A_U$ 是\Term{fibered-category}{纤维范畴} $p^{-1}(U)$。

对 $\mathscr B$ 中的 $f:U\to V$ 以及满足 $p(\lambda)=U,p(\mu)=V$ 的 $\lambda,\mu\in\operatorname{Ob}\mathscr A$，置
\[\operatorname{Hom}_f(\lambda,\mu)=p^{-1}(f),\quad\text{其中 }p:\operatorname{Hom}(\lambda,\mu)\to\operatorname{Hom}(U,V).\]
\begin{SGAresult}{定义}{8.1.1}
设 $\mathscr S$ 是\Term{site}{景}。$\mathscr S$ 上的\Term{locally-cofiltered}{局部余滤过范畴}（或\Term{local-scheme}{局部}左\Term{filtered-category}{滤过范畴}），指 $\mathscr S$ 上的\Term{category}{范畴} $p:\mathscr L\to\mathscr S$，满足：

（L0）对 $\mathscr S$ 中任意的 $f:U\to V$ 及 $\mu\in\operatorname{Ob}\mathscr L_V$，存在 $\lambda\in\operatorname{Ob}\mathscr L_V$，使 $\operatorname{Hom}_f(\lambda,\mu)\ne\varnothing$。

（L1）对任意 $U\in\operatorname{Ob}\mathscr S$ 及 $\mathscr L_U$ 中的\Term{finite}{有限}\Term{object}{对象}族 $(\lambda_i)$，存在 $U$ 的\Term{covering}{覆盖} $f_j:U_j\to U$ 与\Term{object}{对象} $\mu_j\in\operatorname{Ob}\mathscr L_{U_j}$，使对一切 $i,j$，都有 $\operatorname{Hom}_{f_j}(\mu_j,\lambda_i)\ne\varnothing$。

（L2）对 $\mathscr S$ 中任意的 $f:U\to V$ 及位于 $f$ 上方的双\Term{arrow}{箭头} $(\phi_0,\phi_1):\lambda\rightrightarrows\mu$，存在 $V$ 的\Term{covering}{覆盖} $f_j:V_j\to V$ 及靶为 $\lambda$ 的\Term{arrow}{箭头} $\psi_j$，使 $p(\psi_j)\ne f_j$ 且 $\phi_0\psi_j=\phi_1\psi_j$。
\end{SGAresult}
\textbf{8.1.1.1。}当 $\mathscr L$ 在 $\mathscr S^*$ 上为\Term{fibered-category}{纤维范畴}时，此定义可以简化：公理（L0）自动满足，而（L1）、（L2）可表述为：

（L'1）对任意 $U\in\operatorname{Ob}\mathscr S$ 及 $\mathscr L_U$ 中的\Term{finite}{有限}\Term{object}{对象}族 $\lambda_i$，在 $U$ 上\Term{local-scheme}{局部}地存在 $\mu$，它映到所有 $\lambda_i$。

（L'2）对任意 $U\in\operatorname{Ob}\mathscr S$，$\mathscr L_U$ 中的每个双\Term{arrow}{箭头} $(\phi_0,\phi_1):\lambda\rightrightarrows\mu$，在 $U$ 上\Term{local-scheme}{局部}地可由一个靶为 $\lambda$ 的\Term{arrow}{箭头}等化。

\textbf{8.1.2。}设 $\mathscr S$ 是\Term{site}{景}，$\mathscr L$ 是 $\mathscr S$ 上的\Term{category}{范畴}，$\mathscr C$ 是 $\mathscr S$ 上的\Term{stack}{叠}。以 $C$ 表示 $\mathscr C$ 的\Term{global-section}{整体截面}\Term{category}{范畴} $\mathscr Hom_{\mathscr S}^{\mathrm{cart}}(\mathscr S,\mathscr C)$。若 $\mathscr S$ 有\Term{final-object}{终对象} $S$，则 $C$“就是”$\mathscr C_S$。
\begin{SGAresult}{定义}{8.1.2.1}
以 $\mathscr L$ 为指标的 $\mathscr C$ 中\Term{object}{对象}的\Term{local-projective-system}{局部投射系统范畴}，指\Term{category}{范畴} $\mathscr Hom_{\mathscr S}(\mathscr L,\mathscr C)$。
\end{SGAresult}
以 $c$（“相应的常\Term{local-projective-system}{局部投射系统}”）表示复合\Term{functor}{函子}
\[C=\mathscr Hom_{\mathscr S}^{\mathrm{cart}}(\mathscr S,\mathscr C)\hookrightarrow\mathscr Hom_{\mathscr S}(\mathscr S,\mathscr C)\xrightarrow{u\mapsto u\circ p}\operatorname{Hom}_{\mathscr S}(\mathscr L,\mathscr C).
\]
\SourcePage{56}
需要明确对 $\mathscr L$ 的依赖时，改记作 $c_{\mathscr L}$。
\begin{SGAresult}{定义}{8.1.3}
\Term{local-projective-limit}{局部投射极限函子}记作 $\varprojlim_{\lambda\in\mathscr L}$，是\Term{functor}{函子} $c$ 的部分定义的\Term{right-adjoint}{右伴随}\Term{functor}{函子}。
\end{SGAresult}
因此\Term{functor}{函子} $\varprojlim:\mathscr Hom_{\mathscr S}(\mathscr L,\mathscr C)\to C$ 满足
\[\operatorname{Hom}(X,\varprojlim_{\lambda\in\mathscr L}X_\lambda)\simeq\operatorname{Hom}(c(X),(X_\lambda)_{\lambda\in\mathscr L}).\]
\begin{SGAresult}{定义}{8.1.4}
$\mathscr S$ 上的\Term{locally-cofiltered}{局部余滤过范畴}之间的 $\mathscr S$-\Term{functor}{函子} $F:\mathscr L\to\mathscr M$ 称为\Term{cofinal-functor}{共尾函子}，若它满足下述两个条件：

（CF1）对任意 $U\in\operatorname{Ob}\mathscr S$ 及 $\mu\in\operatorname{Ob}\mathscr M_U$，存在 $U$ 的\Term{covering}{覆盖} $f_j:U_j\to U$ 及\Term{object}{对象} $\lambda_j\in\operatorname{Ob}\mathscr L_{U_j}$，使 $\operatorname{Hom}_{f_j}(F(\lambda_j),\mu)\ne\varnothing$。

（CF2）对 $\mathscr S$ 中任意的 $f:U\to V$ 及位于 $f$ 上方的双\Term{arrow}{箭头} $(\phi_0,\phi_1):f(\lambda)\rightrightarrows\mu$，存在 $U$ 的\Term{covering}{覆盖} $f_j:U_j\rightrightarrows U$ 及 $\mathscr L$ 中靶为 $\lambda$ 的\Term{arrow}{箭头} $\psi_j$，使 $p(\psi_j)=f_j$ 且 $\phi_0F(\psi_j)=\phi_1F(\psi_j)$。
\end{SGAresult}
两个\Term{cofinal-functor}{共尾函子}的复合仍为\Term{cofinal-functor}{共尾函子}；\Term{locally-filtered}{局部滤过范畴}的\Term{equivalence}{等价}是\Term{cofinal}{共尾}的（证明留给读者）。
\begin{SGAresult}{命题}{8.1.5}
设 $F:\mathscr L\to\mathscr M$ 是 $\mathscr S$ 上\Term{locally-cofiltered}{局部余滤过范畴}的\Term{cofinal-functor}{共尾函子}，$\mathscr C$ 是 $\mathscr S$ 上的\Term{stack}{叠}，$(X_\mu)_{\mu\in\mathscr M}$ 是以 $\mathscr M$ 为指标的\Term{local-projective-system}{局部投射系统}。则从 $\Gamma\mathscr C$ 到\Term{set}{集合}\Term{category}{范畴}、关于 $X$ 的\Term{functor}{函子}\Term{morphism}{态射}
\[F^*:\operatorname{Hom}(c_{\mathscr M}(X),(X_\mu)_{\mu\in\mathscr M})\longrightarrow\operatorname{Hom}(c_{\mathscr L}(X),(X_{F(\lambda)})_{\lambda\in\mathscr L})\]
是\Term{isomorphism}{同构}。因此
\[F^*:\varprojlim_{\mathscr M}X_\mu\longrightarrow\varprojlim_{\mathscr L}X_{F(\lambda)}\]
是\Term{isomorphism}{同构}，两边同时有定义或同时无定义。
\end{SGAresult}
证明要使用下述引理，其证明留给读者。
\begin{SGAresult}{引理}{8.1.6}
若 $F:\mathscr L\to\mathscr M$ \Term{cofinal}{共尾}，则对 $\mathscr S$ 中任意的 $f:U\to V$ 及位于 $f$ 上方的\Term{arrow}{箭头} $\phi_i:F(\lambda_i')\to\mu$（$i=1,2$），存在 $U$ 的\Term{covering}{覆盖}
\SourcePage{57}
$f_j:U_j\to U$、\Term{object}{对象} $\lambda_j'\in\operatorname{Ob}\mathscr L_{U_j}$ 及\Term{commutative-diagram}{交换图}
\[\begin{tikzcd} &F(\lambda_1)\arrow[dr,"\phi_1"]&\\
F(\lambda_j')\arrow[ur,dashed]\arrow[dr,dashed]&&\mu\\
&F(\lambda_2)\arrow[ur,"\phi_2"']&
\end{tikzcd}\]
在 $\mathscr S$ 中\Term{projection}{投影}为 $(f_j,f)$，即 $U_j\xrightarrow{f_j}U\to V$。
\end{SGAresult}
\begin{proof}[命题 8.1.5 的证明]
对任意 $U\in\operatorname{Ob}\mathscr S$ 和 $\mu\in\mathscr M_U$，由（CF1）存在 $U$ 的\Term{covering}{覆盖} $f_j:U_j\to U$ 及位于 $f_j$ 上方的\Term{arrow}{箭头} $\phi_j:F(\lambda_j)\to\mu$。若
\[\psi=(\psi_\lambda)\in\operatorname{Hom}(c_{\mathscr L}X,(X_{F(\lambda)})_{\lambda\in\mathscr L}),\qquad\psi_\lambda:X|p(\lambda)\to X_{F(\lambda)}\]
是
\[\phi=(\phi_\mu)\in\operatorname{Hom}(c_{\mathscr M}X,(X_\mu)_{\mu\in\mathscr L}),\qquad\phi_\mu:X|p(\mu)\to X_{F(\mu)}\]
的像，则图
\[\begin{tikzcd}&X_{F(\lambda_j)}\arrow[dr,"\phi_j"]&\\
X|U_j\arrow[ur,"\psi_{\lambda_j}"]\arrow[rr,"\phi_\mu|U_j"']&&X_\mu|U_j
\end{tikzcd}\]
交换，所以 $\phi_\mu$ \Term{local-scheme}{局部}地由 $\psi$ 确定；由于 $\mathscr C$ 是\Term{stack}{叠}，就知道 $F^*$ 为\Term{injectivity}{单射}。为证明 $F^*$ \Term{surjectivity}{满射}，从 $\psi$ 出发构造满足 $\psi=F^*\phi$ 的 $\phi$。上述构造给出复合\Term{arrow}{箭头}
\[\phi_{\mu,j}:\phi_j\circ\psi_{\lambda_j}:X|U_j\longrightarrow X_\mu|U_j.\]
\SourcePage{58}
验证这些\Term{arrow}{箭头}可\Term{gluing}{粘合}，从而存在 $\phi_\mu:X|U\to X_\mu$ 使 $\phi_\mu|U_j=\phi_{\mu,j}$。设有\Term{commutative-diagram}{交换图}
\[\begin{tikzcd}&U_{ij}\arrow[dl,"g_i"']\arrow[dr,"g_j"]&\\
U_i\arrow[dr,"f_i"']&&U_j\arrow[dl,"f_j"]\\&U&
\end{tikzcd}\]
为验证 $\phi_{\mu,n}|U_{ij}=\phi_{\mu,i}|U_{ij}$，只需在 $U_{ij}$ 上\Term{local-scheme}{局部}验证。由 8.1.6 和（L0），这使我们可限于在前图上方存在\Term{commutative-diagram}{交换图}
\[\begin{tikzcd}&F(\lambda_{ij})\arrow[dl,"F(a_{ij})"']\arrow[dr,"F(a_{ji})"]&\\
F(\lambda_i)\arrow[dr,"\phi_i"']&&F(\lambda_j)\arrow[dl,"\phi_j"]\\&\mu&
\end{tikzcd}\]
的情形。于是图
\[\begin{tikzcd}&X_{F(\lambda_{ij})}\arrow[dl]\arrow[dr]&\\
X_{F(\lambda_i)}\arrow[dr]&X|U_{ij}\arrow[u]\arrow[l]\arrow[r]\arrow[d]&X_{F(\lambda_j)}\arrow[dl]\\&X_\mu&
\end{tikzcd}\]
\SourcePage{59}
交换。因此存在 $(\phi_\mu)$，使对位于 $f:V\to U$ 上方的每个 $\phi:F(\lambda)\to\mu$，图
\[\begin{tikzcd}&X_{F(\lambda)}\arrow[dr,"\phi"]&\\
X|V\arrow[ur,"\psi_\lambda"]\arrow[rr,"\phi_\mu|V"']&&X_\mu|V
\end{tikzcd}\]
交换。若 $\sigma:\mu\to\mu'$ 是 $\mathscr M$ 的\Term{arrow}{箭头}，则
\[(\sigma|V)(\phi_\mu|V)=(\sigma|V)\phi\cdot\psi_\lambda=(\sigma\phi|V)\cdot\psi_\lambda=\phi_{\mu'}|V,\]
所以 $\phi=(\phi_\mu)$ 是\Term{functor}{函子}\Term{morphism}{态射}，且按要求有 $\psi=F^*\phi$。
\end{proof}
若 $E$ 是\Term{ordered-set}{有序集}，也以 $E$ 表示以 $E$ 为\Term{object}{对象}集的\Term{category}{范畴}，其中当 $i\leq j$ 时 $\operatorname{Hom}(i,j)$ 只有一个元素，否则为空。
\begin{SGAresult}{命题}{8.1.7}
设 $\mathscr L$ 是\Term{site}{景} $\mathscr S$ 上的\Term{locally-cofiltered}{局部余滤过范畴}。存在一个\Term{ordered-set}{有序集} $E$ 及从 $E$ 到 $\mathscr L$ 的\Term{functor}{函子} $F$，满足：

（i）借助 $pF:E\to\mathscr S$ 将 $E$ 看作 $\mathscr S$ 上的\Term{category}{范畴}时，它是\Term{local-scheme}{局部}\Term{cofiltered-category}{余滤过}的。

（ii）\Term{functor}{函子} $F$ \Term{cofinal}{共尾}。
\end{SGAresult}
\begin{proof}
设 $\mathscr L'$ 是 $\mathscr L$ 与\Term{ordered-set}{有序集} $\mathbb Z$ 定义的\Term{category}{范畴}的乘积所成的\Term{locally-cofiltered}{局部余滤过范畴}。\Term{projection}{投影} $\mathscr L'\to\mathscr L$ \Term{cofinal}{共尾}，故只需对 $\mathscr L'$ 证明 (8.1.7)。

若 $X$ 是 $\mathscr L'$ 中的\Term{finite}{有限}\Term{arrow}{箭头}集，以 $X^\circ$ 表示 $X$ 中\Term{arrow}{箭头}的端点集。设 $E$ 是 $\operatorname{Fl}(\mathscr L')$ 的\Term{nonempty}{非空}\Term{finite}{有限}\Term{subset}{子集}所成的\Term{set}{集合}，这些\Term{subset}{子集}满足存在 $\lambda_X\in\operatorname{Ob}\mathscr L'$，使：

a）除 $1_{\lambda_X}$ 外，$X$ 中没有\Term{arrow}{箭头}终止于 $\lambda_X$；

b）若 $\mu\in X^\circ$，则 $1_\mu\in X$ 且 $\operatorname{Hom}(\lambda_X,\mu)\cap X\ne\varnothing$；

c）每个由 $X$ 中\Term{arrow}{箭头}组成的、形如
\[\begin{tikzcd}\lambda_X\arrow[r]\arrow[dr]&\mu\arrow[d]\\&\mu'\end{tikzcd}\]
的图都交换。

用包含关系的反关系在 $E$ 上定义次序。若 $X\in E$ 且 $\mu\in X^\circ$，则从 $\lambda_X$ 到 $\mu$ 的\Term{arrow}{箭头}在 $X$ 中恰有一个。若 $X,Y\in E$
\SourcePage{60}
且 $X\supset Y$，以 $\lambda_{X,Y}$ 表示 $X$ 中从 $\lambda_X$ 到 $\lambda_Y$ 的唯一\Term{arrow}{箭头}。映射 $X\mapsto\lambda_X$ 与 $(X\supset Y)\mapsto\lambda_{X,Y}$ 构成从 $E$ 到 $\mathscr L'$ 的\Term{functor}{函子}。
\end{proof}
\begin{SGAresult}{引理}{8.1.8}
\Term{category}{范畴} $E$ 在 $\mathscr S$ 上\Term{local-scheme}{局部}\Term{cofiltered-category}{余滤过}，而\Term{functor}{函子} $\lambda:E\to\mathscr L'$ \Term{cofinal}{共尾}。
\end{SGAresult}
\begin{proof}依次验证各公理。

（i）公理（L0）。设 $f:V\to U$，$X\in E_U$。存在 $\lambda\in\operatorname{Ob}\mathscr L'_V$ 及 $\phi\in\operatorname{Hom}_f(\lambda,\lambda_X)$。此外，可选择 $\lambda\notin X^\circ$（此处用到 $\mathscr L'=\mathscr L\times\mathbb Z$）。设 $X'=X\cup\{1_\lambda\}\cup Y$，其中 $Y$ 是复合\Term{arrow}{箭头} $\lambda\xrightarrow{\phi}\lambda_X\xrightarrow{\psi}\mu$（$\psi\in X$）的\Term{set}{集合}。于是 $X'\in E$，$\lambda_{X'}=\lambda$，且 $\operatorname{Hom}(X',X)\ne\varnothing$。

（ii）公理（L1）。设 $X_i$ 是 $E_U$ 的\Term{finite}{有限}\Term{object}{对象}族。存在\Term{covering}{覆盖} $f_j:U_j\to U$、\Term{object}{对象} $\lambda_j\in\operatorname{Ob}\mathscr L_{U_j}$ 及\Term{arrow}{箭头} $\phi_{ji}\in\operatorname{Hom}_{f_j}(\lambda_j,\lambda_{X_i})$。若 $x\in X_a^\circ\cap X_b^\circ$，则对每个 $j$，存在\Term{covering}{覆盖} $g_{jk}:U_{jk}\to U_j$ 及\Term{arrow}{箭头} $\psi_{jk}:\lambda_{jk}\to\lambda_j$，使 $p(\psi_{jk})=g_{jk}$，并等化双\Term{arrow}{箭头} $\lambda_j\rightrightarrows x$，其中一个\Term{arrow}{箭头}属于 $X_a$，另一个属于 $X_b$。可安排 $\lambda_{jk}$ 不属于任何 $X_i^\circ$。用 $(f_jg_{jk},\lambda_{jk},\phi_{ji}\psi_{jk})$ 替换 $(f_j,\lambda_j,\phi_{ji})$，并对 $X_i$ 的交中的所有 $x$ 重复此构造。得到一个新系统 $(f_j,\lambda_j,\phi_{ji})$；置
\[X_j=\bigcap_iX_i\cup\{1_{\lambda_j}\}\cup Y_j,\]
其中 $Y_j$ 是复合\Term{arrow}{箭头} $\lambda_j\xrightarrow{\phi_{ji}}\lambda_{X_i}\xrightarrow{\psi}x$（$\psi\in X_i$）的\Term{set}{集合}。于是 $X_j\in E_{V_j}$，且这些 $X_j$ 满足（L1）。

（iii）公理（L3）。由于不存在非平凡的双\Term{arrow}{箭头}，结论显然！

（iv）公理（CF1）。由于\Term{functor}{函子} $\lambda$ 为\Term{surjectivity}{满射}，结论显然。

（v）公理（CF2）。设 $f:U\to V$，并设 $(\phi_0,\phi_1):\lambda_X\rightrightarrows\lambda$ 是位于 $f$ 上方的双\Term{arrow}{箭头}。存在 $U$ 的\Term{covering}{覆盖} $f_j:U_j\to U$ 及\Term{arrow}{箭头} $\psi_j:\lambda_j\to\lambda_X$，使 $p(\psi_j)=f_j$，$\phi_0\psi_j=\phi_1\psi_j$，且 $\lambda_j\notin X^\circ$。设 $X_j=X\cup\{1_{\lambda_j}\}\cup Y_j$，其中 $Y_j$ 是复合\Term{arrow}{箭头} $\lambda_j\xrightarrow{\psi_j}\lambda_X\xrightarrow{\psi}\mu$（$\psi\in X$）的\Term{set}{集合}。于是 $X_j\in E_{V_j}$，且 $F(\lambda_{X_j,X})=\psi_j$ 等化 $(\phi_0,\phi_1)$。
\end{proof}
\textbf{8.1.9.0。}假设 $\mathscr S$ 是单点\Term{topological-space}{拓扑空间}（的\Term{nonempty}{非空}\Term{open-subset}{开集}\Term{site}{景}）。$\mathscr S$ 上的\Term{locally-cofiltered}{局部余滤过范畴}就只是左\Term{filtered-category}{滤过范畴}（即其 $\mathscr S^\circ$ 在第 I 讲 2.7 的意义下\Term{filtered}{滤过}）；\Term{cofinal-functor}{共尾函子}则对应于左\Term{filtered-category}{滤过范畴}的\Term{cofinal-functor}{共尾函子}（第 I 讲 8）。命题 8.1.7 可重述为：
\begin{SGAresult}{命题}{8.1.9}
对每个余\Term{filtered-category}{滤过范畴} $\mathscr L$，存在一个\Term{cofiltered-category}{余滤过}\Term{ordered-set}{有序集} $E$ 及从 $E$ 到 $\mathscr L$ 的\Term{cofinal-functor}{共尾函子}。
\end{SGAresult}
\SourcePage{61}
在这一特殊情形，8.1.7 的证明实质上归结为第 I 讲 8 给出的 (8.1.9) 的证明。

\textbf{8.1.10。}设 $\mathscr T\to\mathscr S$ 是\Term{site}{景}的\Term{morphism}{态射}，$p:\mathscr L\to\mathscr S$ 是 $\mathscr S$ 上的\Term{locally-cofiltered}{局部余滤过范畴}。以 $\mathscr L^*$ 表示下述\Term{category}{范畴}：

（i）$\mathscr L^*$ 的\Term{object}{对象}是四元组 $(V,U,\phi,\lambda)$，其中 $V\in\operatorname{Ob}\mathscr T$，$U\in\operatorname{Ob}\mathscr S$，$\phi\in\operatorname{Hom}(V,q^*U)$，$\lambda\in\operatorname{Ob}\mathscr L_U$。

（ii）\Term{arrow}{箭头} $f:(V,U,\phi,\lambda)\to(V',\phi',U',\lambda')$ 是三元组 $(f_1,f_2,f_3)$，其中 $f_1\in\operatorname{Hom}(V,V')$，$f_2\in\operatorname{Hom}(U,U')$，$f_3\in\operatorname{Hom}(\lambda,\lambda')$，且 $q^*f_2\circ\phi=\phi'\circ f_1$，$p(f_3)=f_2$。

\Term{functor}{函子} $f\mapsto f_1$ 使 $\mathscr L^*$ 成为 $\mathscr T$ 上的\Term{category}{范畴}。
\begin{SGAresult}{引理}{8.1.11}
\Term{category}{范畴} $\mathscr L^*$ 在 $\mathscr T$ 上\Term{local-scheme}{局部}\Term{cofiltered-category}{余滤过}。
\end{SGAresult}
\begin{proof}
公理（L0）显然成立。验证（L1）：设有\Term{finite}{有限}族 $L_i=(V,U_i,\phi_i,\lambda_i)$。各 $\phi_i$ 定义从 $V$ 到\Term{sheaf}{层} $q^*a\prod U_i=a\prod q^*U_i$ 的\Term{morphism}{态射}，故存在\Term{commutative-diagram}{交换图}
\[\begin{tikzcd}V_j\arrow[r,"f_j"]\arrow[d,"\psi_j"']&V\arrow[d,"\phi_i"]\\q^*U_j\arrow[r,"q^*(P_{ij})"']&q^*U_i\end{tikzcd}\]
使 $f_j$ \Term{covering}{覆盖} $V$。

取 $\lambda_{ji}$ 使 $\operatorname{Hom}_{p_{ij}}(\lambda_{ji},\lambda_i)\ne\varnothing$。由（L1），在适当\Term{refinement}{加细} $U_j$（及 $V_j$）后，可取不依赖于 $i$ 的 $\lambda_{ji}=\lambda_j$；此时 $L_j=(V_j,U_j,\psi_j,\lambda_j)$ 满足（L1）。

验证（L2）。设 $(f_1,f_2):L^1\rightrightarrows L^2$，其中 $L^i=(V^i,U^i,\phi^i,\lambda^i)$，$f_i=(f,g_i,h_i)$。各 $\phi_i$ 定义从 $V$ 到\Term{sheaf}{层} $q^*a\operatorname{Ker}(U^1\rightrightarrows U^2)=a\operatorname{Ker}(q^*U^1\rightrightarrows q^*U^2)$ 的\Term{morphism}{态射}，所以存在\Term{commutative-diagram}{交换图}
\[\begin{tikzcd}V_j\arrow[r,"f_j"]\arrow[d,"\psi_j"']&V^1\arrow[r]\arrow[d,"\phi^1"]&V^2\arrow[d,"\phi^2"]\\
q^*U_j\arrow[r,"p_j"']&q^*U^1\arrow[r,shift left]\arrow[r,shift right]&q^*U^2
\end{tikzcd}\quad(g_1p_j=g_2p_j),\]
其中 $f_j$ \Term{covering}{覆盖} $V$。取 $\chi_j:\lambda^j\to\lambda^1$，使 $p(\chi_j)=p_j$。将（L2）应用于 $(\lambda_1\chi_j,h_2\chi_j)$，适当\Term{refinement}{加细} $U_j$（及 $V_j$）后，可安排 $h_1\chi_j=h_2\chi_j$；此时 $L_j=(V_j,U_j,\psi_j,\lambda_j)$ 满足（L2）。
\end{proof}
\SourcePage{62}
\begin{SGAresult}{命题}{8.1.12}
设 $t^*:\mathscr S^*\to\mathscr S$ 是 $\mathscr S$ 上的\Term{locally-cofiltered}{局部余滤过范畴}，并在 $\mathscr S^*$ 上赋予由 $\mathscr S$ 的\Term{topology}{拓扑}经 $t^*$ 诱导的\Term{topology}{拓扑}。则 $t^*$ 是\Term{site}{景}的\Term{equivalence}{等价} $t:\mathscr S\to\mathscr S^*$。
\end{SGAresult}
\begin{proof}
由构造，$t_*:F\mapsto F\circ t^*$ 把 $\mathscr S$ 上的\Term{sheaf}{层}变为 $\mathscr S^*$ 上的\Term{sheaf}{层}。由（L0）以及应用于空族的（L1），$\mathscr S$ 的每个“足够小”的\Term{object}{对象}都在 $t^*$ 的像中（即 $t^*$ 的像是 $\mathscr S$ 中的\Term{covering}{覆盖}筛）。所以由“\Term{comparison-lemma}{比较引理}”（第 III 讲 4），将 $t^*(\operatorname{Ob}\mathscr S^*)$ 定义的\Term{full-subcategory}{满子范畴}包含到 $\mathscr S$ 中是\Term{site}{景}的\Term{equivalence}{等价}。因而可以假设 $t^*$ 在\Term{object}{对象}上为\Term{surjectivity}{满射}。

设 $F$ 是 $\mathscr S^*$ 上的\Term{sheaf}{层}。

（i）设 $f_j:U_j\to U$ 是 $\mathscr S$ 中的\Term{covering}{覆盖}，$\lambda^1,\lambda^2\in\mathscr S_U^*$，且对 $i=1,2$，有 $f_j^i\in\operatorname{Hom}_{f_j}(\mu_j,\lambda^i)$。则 $F(\lambda^1)$ 与 $F(\lambda^2)$ 在 $\prod_jF(\mu_j)$ 中的像相同。

事实上，设 $x=(x_j)$ 属于 $F(\lambda^1)$ 的像。为了使 $x$ 属于 $F(\lambda^2)$ 的像，只需对每个\Term{commutative-diagram}{交换图}
\[\begin{tikzcd}\mu\arrow[r]\arrow[d]&\mu_j\arrow[d]\\\mu_k\arrow[r]&2\end{tikzcd}\]
都有 $x_j,x_k$ 在 $F(\mu)$ 中的像相同。由于 $F$ 是\Term{sheaf}{层}，只需将 $\mu$ 替换为其一个\Term{covering}{覆盖}的各\Term{object}{对象}后验证；由（L2），这使我们可假设将前图中的 $\lambda^2$ 替换为 $\lambda^1$ 所得的类似图也交换，此时断言显然。

（ii）设 $f:U\to V$ 在 $\mathscr S^*$ 中，$\lambda\in\operatorname{Ob}\mathscr S_U^*$，$\mu\in\operatorname{Ob}\mathscr S_V^*$。则存在\Term{covering}{覆盖} $f_j:U_j\to U$、\Term{object}{对象} $\mu'\in\operatorname{Ob}\mathscr S_U^*$ 及\Term{arrow}{箭头}
\[\begin{tikzcd}&\lambda&\\\lambda_j\arrow[ur,"\phi_j"]\arrow[dr,"\psi_j"']&&\mu\\&\mu'\arrow[ur,"\psi"']&\end{tikzcd}\]
位于 $U_j\xrightarrow{f_j}U\xrightarrow{f}V$ 上方。
\SourcePage{63}
由（i），存在唯一\Term{arrow}{箭头} $\overline f$，使下图交换：
\[\begin{tikzcd}F(\mu)\arrow[r,"\overline f"]\arrow[d,"\psi"']&F(\lambda)\arrow[d,"\phi_j"]\\
F(\mu')\arrow[r,"\psi_j"']&F(\lambda_j)\end{tikzcd}\]
若将 $(f_j)$ 替换为由 $f_k:U_k\to U_{j(k)}$ 给出的更细\Term{covering}{覆盖}，将 $\lambda_j$ 替换为带有 $\phi_k\in\operatorname{Hom}_{f_k}(\lambda_k,\lambda_{j(k)})$ 的 $\lambda_k$，并将 $\phi_j$（$\psi_j$）替换为 $\phi_{j(k)}\circ\phi_k$（$\psi_{j(k)}\circ\phi_k$），则\Term{arrow}{箭头} $\overline f$ 不变。

由此推出 $\overline f$ 只依赖于 $f$，留给读者证明。

（iii）证明 $\overline f$ 关于 $f$ \Term{functor}{函子}性。设 $W\xrightarrow{h}V\xrightarrow{g}U$ 在 $\mathscr S$ 中，且 $\lambda^0,\mu^0,\nu^0$ 分别属于 $\operatorname{Ob}\mathscr S_U^*,\operatorname{Ob}\mathscr S_V^*,\operatorname{Ob}\mathscr S_W^*$。存在\Term{commutative-diagram}{交换图}
\[\begin{tikzcd}[column sep=large,row sep=large]
\nu^0&&\mu^0&\lambda^0\\
&\nu^1\arrow[ur]&&\mu^1\arrow[u]\\
&\nu^2\arrow[urr]&&\\
\nu_{jk}\arrow[uuu,"e_{jk}^0"]\arrow[uur,"e_{jk}^1"]\arrow[ur,"e_{jk}^2"]&&&\mu_j\arrow[uuul,"f_j^0"']\arrow[uu,"f_j^1"']\arrow[from=lll]
\end{tikzcd}\]
满足：

a）$t^*(\nu^i),t^*(\mu^i),t^*(e_{jk}^i),t^*(f_j^i)$ 不依赖于 $i$；

b）固定 $i$ 后，各 $e_{jk}^i$ \Term{covering}{覆盖} $W$，各 $f_j^i$ \Term{covering}{覆盖} $V$。

先用（L0）构造 $\mu^1,\nu^1,\nu^2$，再按（ii）构造 $\mu_j$。设 $W'_{jk}$ 是 $W$ 的一个\Term{covering}{覆盖}，它按图所要求的方式映到\Term{covering}{覆盖} $V_j=t^*(\mu_j)$。\Term{refinement}{加细} $W'_{jk}$ 并应用（L0）与（L1），得到一个未必交换、但形状符合要求且满足 a）、b）的图。再次\Term{refinement}{加细}并应用（L2），就使它成为\Term{commutative-diagram}{交换图}。
\SourcePage{64}
于是只需观察下列\Term{commutative-diagram}{交换图}：
\[\begin{tikzcd}[column sep=large,row sep=large]
F(\lambda^0)\arrow[r,"\overline g"]&F(\mu^0)\arrow[rr,"\overline h"]\arrow[dr]&&F(\nu^0)\arrow[ddd]\\
&&F(\nu^1)\arrow[ddr]&\\
F(\mu^1)\arrow[r]&\prod F(\mu_j)\arrow[drr]&&\\
F(\nu^2)\arrow[rrr]&&&\prod F(\nu_{jk}).
\end{tikzcd}\]
（iv）$\overline g$ 显然关于 $F$ \Term{functor}{函子}性：因此，$\mathscr S^*$ 上的每个\Term{sheaf}{层} $F$ 都是 $\mathscr S$ 上一个（唯一确定的）\Term{presheaf}{预层} $F_0$ 在 $t$ 下的\Term{direct-image}{正像}，而\Term{sheaf}{层}的每个\Term{morphism}{态射} $f:F\to G$ 都是唯一一个\Term{presheaf}{预层}\Term{morphism}{态射} $f_0:F_0\to G_0$ 的像。

由此已知 $t_*$ \Term{fully-faithful}{全忠实}；尚须证明 $F_0$ 是\Term{sheaf}{层}。若 $f_j:U_j\to U$ 是\Term{covering}{覆盖}，则此\Term{covering}{覆盖}是 $\mathscr S^*$ 中某个\Term{covering}{覆盖}的像，所以 $F_0(U)$ \Term{injectivity}{单射}到 $\prod_jF_0(U_j)$ 中。若 $x_j\in F_0(U_j)$ 可\Term{gluing}{粘合}，那么它们在 $\mathscr S^*$ 中当然也可\Term{gluing}{粘合}，从而来自 $F_0(U)$ 的某个元素。至此完成 (8.1.12) 的证明。
\end{proof}
\textbf{8.1.13.0。}设 $q:\mathscr T\to\mathscr S$、$p:\mathscr L\to\mathscr S$ 及 $\mathscr L^*$ 如 (8.1.10)，$\mathscr C$ 是 $\mathscr T$ 上的\Term{stack}{叠}。$\mathscr S$ 上的\Term{stack}{叠} $q_*\mathscr C$ 是由 $(q_*\mathscr C)_U=\mathscr C_{q^*U}$“定义”的\Term{stack}{叠}。

若 $(X_\lambda)_{\lambda\in\mathscr L}$ 是以 $\mathscr L$ 为指标、取值于 $\mathscr C$ 的\Term{local-projective-system}{局部投射系统}，则按公式
\[X_L=\phi^*X_\lambda\qquad\text{对 }L=(V,U,\phi,\lambda)\]
定义 $(X_L)_{L\in\mathscr L^*}$。下述命题留给读者验证。
\begin{SGAresult}{命题}{8.1.13.1}
沿用上述记号，有
\[\varprojlim_{\lambda\in\mathscr L}X_\lambda\quad\varprojlim_{L\in\mathscr L^*}X_L,\]
两边同时有定义或同时无定义。
\end{SGAresult}
设 $\mathscr S$ 是\Term{site}{景}，$p^*:\mathscr S^*\to\mathscr S$ 是 $\mathscr S$ 上的\Term{locally-filtered}{局部滤过范畴}。若 $\mathscr C$ 是 $\mathscr S$ 上的\Term{stack}{叠}，则 $p_*\mathscr C$ 是 $\mathscr S^*$ 上的\Term{stack}{叠}（8.1.13.0）；此外，恒等\Term{functor}{函子} $\mathscr S^*\to\mathscr S^*$ 是赋予\Term{induced-topology}{诱导拓扑}的 $\mathscr S^*$ 上的\Term{locally-filtered}{局部滤过范畴}，而每个以 $\mathscr S^*$ 为指标的\Term{local-inductive-system}{局部归纳系统} $(X_\lambda)$ 都在 $\mathscr S^*$ 上定义一个以 $\mathscr S^*$ 为指标的\Term{local-inductive-system}{局部归纳系统}。

容易验证：
\begin{SGAresult}{命题}{8.1.14}
沿用上述记号，有
\[\varprojlim_{\lambda\in\mathscr S^*/\mathscr S}X_\lambda=\varprojlim_{\lambda\in\mathscr S^*/\mathscr S^*}X_\lambda,\]
\SourcePage{65}
两边同时有定义或同时无定义。
\end{SGAresult}
\textbf{8.2。层\Term{category}{范畴}中的\Term{local-inductive-limit}{局部归纳极限}。}

\textbf{8.2.0。}设 $p:\mathscr T\to\mathscr S$ 是\Term{site}{景}的\Term{morphism}{态射}。设 $\mathscr C$ 是 $\mathscr S$ 上如下的\Term{stack}{叠}：

（i）$\mathscr C$ 的\Term{object}{对象}是二元组 $(U,F)$，其中 $U$ 是 $\mathscr S$ 的\Term{object}{对象}，$F$ 是 $p^*U$ 上的\Term{sheaf}{层}。

（ii）\Term{arrow}{箭头} $f:(U,F)\to(V,G)$ 是由一个\Term{arrow}{箭头} $f_0:U\to V$ 及一个 $\mathscr T/U$-\Term{sheaf}{层}\Term{morphism}{态射} $f_1:p^*(f_0)^*G\to F$ 所成的二元组。

（iii）从 $\mathscr C$ 到 $\mathscr S$ 的\Term{functor}{函子}由 $f\mapsto f_0$ 给出。

应注意，$\mathscr C_U$ 是 $p^*U$ 上的\Term{sheaf}{层}\Term{category}{范畴}的\Term{opposite-category}{反范畴}。

设 $\mathscr L$ 是 $\mathscr S$ 上的\Term{local-scheme}{局部}\Term{cofiltered-category}{余滤过}（左\Term{filtered}{滤过}）\Term{category}{范畴}。
\begin{SGAresult}{定义}{8.2.1}
沿用上述记号，以 $\mathscr L$ 为指标的 $\mathscr T$ 上的\Term{sheaf}{层}的\Term{local-inductive-system}{局部归纳系统范畴}，指\Term{category}{范畴} $\operatorname{Hom}_{\mathscr S}(\mathscr L,\mathscr C)^\circ$。
\end{SGAresult}
8.1.3 的“\Term{projective-limit}{投射极限}”\Term{functor}{函子}在这里称为“\Term{local-inductive-limit}{局部归纳极限}”\Term{functor}{函子}。它是从 $\operatorname{Hom}_{\mathscr S}(\mathscr L,\mathscr C)^\circ$ 到 $(\Gamma\mathscr C)^\circ=\mathscr T^{\sim}$ 的\Term{functor}{函子}。
\begin{SGAresult}{定理}{8.2.2}
从以 $\mathscr L$ 为指标的 $\mathscr T$ 上的\Term{sheaf}{层}的\Term{inductive-system}{归纳系统}\Term{category}{范畴}到 $\mathscr T$ 上的\Term{sheaf}{层}\Term{category}{范畴}的“\Term{local-inductive-limit}{局部归纳极限}”\Term{functor}{函子}，处处有定义，与\Term{finite}{有限}\Term{projective-limit}{投射极限}交换，也与任意\Term{inductive-limit}{归纳极限}交换。
\end{SGAresult}
\begin{proof}
由命题 8.1.13、8.1.14 及 8.1.12，可归约到 $\mathscr S=\mathscr T=\mathscr L$ 的情形。此时\Term{local-inductive-system}{局部归纳系统}就是对每个 $U\in\operatorname{Ob}\mathscr S$ 给定 $\mathscr S/U$ 上的\Term{sheaf}{层} $F_U$，并对 $\mathscr S$ 中每个\Term{arrow}{箭头} $f:V\to U$ 给定关于 $f$ \Term{functor}{函子}性的\Term{arrow}{箭头} $f^*F_U\to F_V$。

对于这样的\Term{inductive-system}{归纳系统}，\Term{inductive-limit}{归纳极限}按下述方式计算。
\end{proof}
\begin{SGAresult}{引理}{8.2.3}
设 $U\mapsto F_U$ 是 $\mathscr S$ 各\Term{object}{对象}上的\Term{presheaf}{预层}\Term{category}{范畴}的一个\Term{section}{截面}。有
\[a(U\mapsto\Gamma(U,F_U))\simeq\varprojlim_U aF_U.\]
\end{SGAresult}
\begin{proof}
根据定义，右边表示的\Term{functor}{函子}，将 $\mathscr S$ 上的每个\Term{sheaf}{层} $F$ 映到由\Term{arrow}{箭头} $\phi_U:F_U\to F|U$ 组成的相容系统的\Term{set}{集合}。进而，一个\Term{arrow}{箭头} $\phi_U$ 就是由\Term{arrow}{箭头} $\phi_g$ 组成的相容系统，其中每个 $g:V\to U$ 对应一个 $\phi_g:F_U(V)\to F(V)$。对 $\mathscr S$ 的恒等\Term{functor}{函子}应用 8.1.11，得到 $\mathscr S$ 上\Term{local-scheme}{局部}\Term{filtered-category}{滤过的} $\mathscr S^*$。已经看到
\[\varinjlim_{g\in\mathscr S^*}F_U(V_g)^{\sim}\simeq\varinjlim_UaF_U,\]
\SourcePage{66}
其中 $V_g$ 是 $g$ 的源，${}^{\sim}$ 是“所生成的常\Term{sheaf}{层}”\Term{functor}{函子}。从 $\mathscr S$ 到 $\mathscr S^*$ 的\Term{functor}{函子} $U\mapsto1_U$ \Term{cofinal}{共尾}，所以由 8.1.5 又有
\[\varinjlim_UF_U(U)^{\sim}\simeq\varinjlim_UaF_U.\]
回到定义，最后得到
\[\varinjlim_UF_U(U)^{\sim}=a(U\mapsto\Gamma(U,F_U)).\]
\end{proof}
这表明，在已归约到的情形，\Term{functor}{函子} $\varinjlim$ 处处有定义，并可作为“由\Term{presheaf}{预层}生成的\Term{sheaf}{层}”\Term{functor}{函子}与\Term{functor}{函子} $(F_U)\mapsto(U\mapsto\Gamma(U,F_U))$ 的复合来计算。这两个\Term{functor}{函子}都与\Term{finite}{有限}\Term{projective-limit}{投射极限}交换，而\Term{functor}{函子} $\varinjlim$ 由其作为\Term{adjunction}{伴随}\Term{functor}{函子}的定义，还与任意\Term{inductive-limit}{归纳极限}交换。这完成了定理 8.2.2 的证明。

引理 8.2.3 提供了一个系统方法：从“\Term{presheaf}{预层}的\Term{associated-sheaf}{伴随层}”\Term{functor}{函子}的类似性质，证明\Term{local-inductive-limit}{局部归纳极限}\Term{functor}{函子}的性质。由此得到：
\begin{SGAresult}{命题}{8.2.4}
设 $\mathscr T,\mathscr S,\mathscr L$ 如前，$A_\lambda$ 是 $\mathscr T$ 上环\Term{sheaf}{层}的\Term{local-inductive-system}{局部归纳系统}，$L_\lambda$（相应地 $M_\lambda$）是 $A_\lambda$ 上右\Term{module}{模}\Term{sheaf}{层}（相应地左\Term{module}{模}\Term{sheaf}{层}）的\Term{local-inductive-system}{局部归纳系统}，三者均以 $\mathscr L$ 为指标。有
\[\varinjlim(L_\lambda\otimes_{A_\lambda}M_\lambda)=\varinjlim L_\lambda\otimes_{\varinjlim A_\lambda}\varinjlim M_\lambda.\]
\end{SGAresult}
（归约到 $\mathscr T=\mathscr S=\mathscr L$，并注意两边都等于由\Term{presheaf}{预层} $L_U(U)\otimes_{A_U(U)}M_U(U)$ 生成的\Term{sheaf}{层}。）
\begin{SGAresult}{命题}{8.2.5}
设 $\mathscr T_1\xrightarrow{q}\mathscr T_2\to\mathscr S$ 是两个\Term{site}{景}的\Term{morphism}{态射}，$(F_\lambda)$ 是 $\mathscr T_2$ 上以 $\mathscr S$ 上的\Term{locally-cofiltered}{局部余滤过范畴} $\mathscr L$ 为指标的\Term{sheaf}{层}的\Term{local-inductive-system}{局部归纳系统}。有
\[q^*\varinjlim_{\lambda\in\mathscr L}F_\lambda\simeq\varinjlim_{\lambda\in\mathscr L}q^*F_\lambda.\]
\end{SGAresult}
\begin{proof}这立即由 $q^*$ 有\Term{right-adjoint}{右伴随} $q_*$ 得到。\end{proof}
设 $\mathscr S,\mathscr T$ 是两个\Term{finite}{有限}\Term{projective-limit}{投射极限}均\Term{representable-functor}{可表示}的\Term{site}{景}，$p:\mathscr T\to\mathscr S$ 是\Term{site}{景}的\Term{morphism}{态射}，使 $p^*:\mathscr S\to\mathscr T$ 与\Term{finite}{有限}\Term{projective-limit}{投射极限}交换；设 $\mathscr L$ 是 $\mathscr T$ 上以三元组 $(V,U,\phi)$ 为\Term{object}{对象}的\Term{locally-cofiltered}{局部余滤过范畴}，其中 $V\in\operatorname{Ob}\mathscr T$，$U\in\operatorname{Ob}\mathscr S$，$\phi\in\operatorname{Hom}(V,p^*U)$。它是对 $1_{\mathscr S}:\mathscr S\to\mathscr S$ 应用 8.1.11 得到的。
\begin{SGAresult}{命题}{8.2.6}
对 $\mathscr T$ 上的每个\Term{sheaf}{层} $F$，有
\[F=\varinjlim_{\lambda\in\mathscr L}\phi^*\phi_*F,\]
\SourcePage{67}
这里滥用记号，以 $\phi$ 表示由 $p$ 诱导的从 $\mathscr T/V$ 到 $\mathscr S/V$ 的\Term{morphism}{态射}。
\end{SGAresult}
\begin{proof}
设 $\phi^0$ 是\Term{presheaf}{预层}意义下的\Term{inverse-image}{逆像}\Term{functor}{函子}；由 8.2.3 得到
\[\varinjlim_{\lambda\in\mathscr L}\phi^{0*}\phi_*F
=\varinjlim_{\lambda\in\mathscr L}a\phi^0\phi_*F
=\varinjlim_{\lambda\in\mathscr L}F(V)^{\sim}
=\varinjlim_VF(V)^{\sim}=F.\]
\end{proof}
\begin{SGAresult}{命题}{8.2.7}
设 $p:T\to S$ 是\Term{topos}{拓扑斯}\Term{morphism}{态射}，$\mathcal O_S$ 是 $S$ 上的环\Term{sheaf}{层}，$\mathcal O_T$ 是它在 $T$ 上的\Term{inverse-image}{逆像}。右 $\mathcal O_T$-\Term{module}{模} $M$ 在 $\mathcal O_T$ 上\Term{flat-module}{平坦}的充要条件是：从左 $\mathcal O_S$-\Term{module}{模}\Term{category}{范畴}到 $T$ 上\Term{abelian-sheaf}{阿贝尔层}\Term{category}{范畴}的\Term{functor}{函子}
\begin{equation}N\longmapsto M\otimes_{\mathcal O_T}p^*N\tag{8.2.7.1}\end{equation}
\Term{exact}{正合}。
\end{SGAresult}
\begin{proof}
必要性显然；因此假设\Term{functor}{函子} (8.2.7.1) \Term{exact}{正合}，证明 $M$ \Term{flat}{平坦}。可以假设 $p$ 由一个满足 2.6 中假设的\Term{site}{景}\Term{morphism}{态射} $p:\mathscr T\to\mathscr S$ 定义。设 $V\in\operatorname{Ob}\mathscr T$，$U\in\operatorname{Ob}\mathscr S$，$\phi:V\to p^*U$。仍以 $\phi$ 表示诱导的\Term{site}{景}\Term{morphism}{态射} $\phi:\mathscr T/V\to\mathscr S/U$。
\end{proof}
\begin{SGAresult}{引理}{8.2.8}
从 $\mathscr S/U$ 上的 $\mathcal O_S|U$-\Term{module}{模}\Term{category}{范畴}到 $\mathscr T/V$ 上的 $\mathcal O_T|V$-\Term{module}{模}\Term{category}{范畴}的\Term{functor}{函子} $N\mapsto M|V\otimes_{\mathcal O_V}\phi^*N$ 是\Term{exact}{正合}的。
\end{SGAresult}
\begin{proof}
可归约到 $V=p^*U$。设 $j,j'$ 是“包含\Term{morphism}{态射}”$j:\mathscr S/U\to\mathscr S$ 与 $j':\mathscr T/p^*U\to\mathscr T$。有
\[j'_!(M|V\otimes_{\mathcal O_V}\phi^*N)=M\otimes_{\mathcal O_T}p^*j_!N.\]
由于 $j_!$ 与 $j'_!$ \Term{exact}{正合}且\Term{faithful}{忠实}，即得 8.2.8。
\end{proof}
\begin{proof}[命题 8.2.7 证明续]
设 $Q$ 是 $\mathcal O_T$-\Term{module}{模}。有（8.2.6）
\[Q=\varinjlim_{\phi:V\to p^*U}\phi^*\phi_*Q|V,\]
从而（8.2.4）
\[M\otimes_{\mathcal O_T}Q=\varinjlim_{\phi:V\to p^*U}M|V\otimes_{\mathcal O_{S_V}}\phi^*\phi_*Q|V.\]
由 8.2.8，\Term{functor}{函子} $Q\mapsto M|V\otimes_{\mathcal O_V}\phi^*\phi_*Q|V$ \Term{left-exact}{左正合}，所以 $M\otimes_{\mathcal O_T}Q$ 关于 $Q$ \Term{left-exact}{左正合}，因而\Term{exact}{正合}。
\end{proof}
\begin{SGAresult}{推论}{8.2.9}
设 $f:(T,\mathcal O_T)\to(S,\mathcal O_S)$ 是\Term{ringed-topos}{赋环拓扑斯}\Term{morphism}{态射}。那么，\Term{flat}{平坦}（右）\Term{module}{模}\Term{sheaf}{层}在 $f$ 下的\Term{inverse-image}{逆像}仍是\Term{flat-module}{平坦模}\Term{sheaf}{层}。
\end{SGAresult}
\SourcePage{68}
\begin{proof}
通过 $(T,f^*\mathcal O_S)$ \Term{resolution}{分解} $f$，可归约到 $\mathcal O_T=f^*\mathcal O_S$ 的情形。设 $M$ 是 $S$ 上的\Term{flat-module}{平坦模}\Term{sheaf}{层}。由 8.2.7，要验证 $f^*M$ \Term{flat}{平坦}，只需验证\Term{functor}{函子}
\[N\longmapsto f^*M\otimes_{\mathcal O_T}f^*N=f^*(M\otimes_{\mathcal O_S}N)\]
的\Term{exactness}{正合性}；由于 $f^*$ \Term{exact}{正合}，结论成立。
\end{proof}
\begin{SGAresult}{引理}{8.2.10}
设 $(S,\mathcal O_S)$ 是\Term{ringed-topos}{赋环拓扑斯}，$F$ 是\Term{locally-finitely-presented}{局部有限表示}的左\Term{module}{模}\Term{sheaf}{层}，$G$ 是\Term{bimodule}{双模}\Term{sheaf}{层}，$H$ 是\Term{flat}{平坦}左\Term{module}{模}\Term{sheaf}{层}。则\Term{canonical}{典范}\Term{arrow}{箭头}
\[\mathscr Hom(F,G)\otimes H\longrightarrow\mathscr Hom(F,G\otimes H)\]
是\Term{isomorphism}{同构}。
\end{SGAresult}
\begin{proof}
问题在 $S$ 上是\Term{local-scheme}{局部}的，故可假设 $F$ 有\Term{finitely-presented}{有限表示}
\[F_1\longrightarrow F_0\longrightarrow F\longrightarrow0.\]
由于 $H$ \Term{flat}{平坦}，下图的第一行\Term{exact}{正合}：
\[\begin{tikzcd}[column sep=small]
0\arrow[r]&\mathscr Hom(F,G)\otimes H\arrow[r]\arrow[d]&\mathscr Hom(F_1,G)\otimes H\arrow[r]\arrow[d,"\simeq"]&\mathscr Hom(F_2,G)\otimes H\arrow[d,"\simeq"]\\
0\arrow[r]&\mathscr Hom(F,G\otimes H)\arrow[r]&\mathscr Hom(F_1,G\otimes H)\arrow[r]&\mathscr Hom(F_2,G\otimes H)
\end{tikzcd}\]
结论随之成立。令 $G=A$，由 8.2.10 得到：
\end{proof}
\begin{SGAresult}{引理}{8.2.11}
从\Term{locally-finitely-presented}{局部有限表示}\Term{module}{模}\Term{sheaf}{层}到\Term{flat-module}{平坦模}\Term{sheaf}{层}的每个\Term{morphism}{态射}，在 $S$ 上\Term{local-scheme}{局部}地都经过一个自由\Term{sheaf}{层} $\mathcal O_S^n$（$n\in\mathbb N$）\Term{resolution}{分解}。
\end{SGAresult}
\begin{SGAresult}{定理}{8.2.12}
（D. Lazard。）
\Term{ringed}{赋环}\Term{site}{景} $(\mathscr S,\mathcal O_S)$ 上的\Term{module}{模}\Term{sheaf}{层} $M$ \Term{flat}{平坦}的充要条件是：它是\Term{finite-type}{有限型}\Term{free-module}{自由模}的\Term{local-inductive-limit}{局部归纳极限}。
\end{SGAresult}
\begin{proof}
充分性由 8.2.2、8.2.4 得到。

对每个 $U\in\operatorname{Ob}\mathscr S$，以 $\mathscr L_0(U)$（相应地 $\mathscr L_1(U)$）表示 $U$ 上附带一个到 $M|U$ 的线性映射的\Term{finite-type}{有限型}\Term{free-module}{自由模}\Term{sheaf}{层}（相应地，\Term{finitely-presented}{有限表示}\Term{module}{模}\Term{sheaf}{层}）的\Term{category}{范畴}。对 $\mathscr S$ 中每个\Term{arrow}{箭头} $f:V\to U$，以 $f^*$ 表示从 $\mathscr L_i(U)$ 到 $\mathscr L_i(V)$ 的\Term{restriction}{限制}\Term{functor}{函子}。这些\Term{functor}{函子}定义 $\mathscr S$ 上的\Term{fibered-category}{纤维范畴} $\mathscr L_i$，其\Term{fiber}{纤维}是\Term{category}{范畴} $\mathscr L_i(U)$ 的\Term{opposite-category}{反范畴}。

读者可验证，\Term{category}{范畴} $\mathscr L_1$ 在 $\mathscr S$ 上\Term{local-scheme}{局部}\Term{filtered}{滤过}，且
\[M=\varinjlim_{(N,f)\in\mathscr L_1}N.\]
\SourcePage{69}
若 $M$ \Term{flat}{平坦}，则 $\mathscr L_0$ 到 $\mathscr L_1$ 的包含\Term{fully-faithful}{全忠实}，且由 8.2.11 满足 8.1.4 的（CF1）。于是它自动满足（CF2），而 $\mathscr L_0$ \Term{local-scheme}{局部}\Term{filtered}{滤过}。由 8.1.5，有
\[M=\varinjlim_{(N,f)\in\mathscr L_2}N,\]
从而得到断言 8.2.12。
\end{proof}
\textbf{参考文献}
\begin{enumerate}
\item[{[1]}] M. Artin 与 B. Mazure：《\Term{etale}{平展}\Term{homotopy}{同伦}论》（Etale Homotopy Theory），Lecture Notes，编号未列，Springer Verlar。
\item[{[2]}] Blum 与 Herrera：将在 Inventiones 发表的文章。
\item[{[3]}] H. Cartan 与 S. Eilenberg：《\Term{homology}{同调}\Term{algebra}{代数}》（Homological Algebra）。
\item[{[4]}] P. Deligne：《Hodge 理论》（Théorie de Hodge），I.H.E.S. 出版物。
\item[{[5]}] P. Gabriel 与 N. Popescu：CRAS。
\item[{[6]}] P. Gabriel 与 M. Zisman：《\Term{homotopy}{同伦}论与分式演算》（Homotopie Theory and Calculus of Fraction），Ergebnisse der Mathematik，第 35 卷。
\item[{[7]}] R. Godement：《\Term{sheaf}{层}论》（Théorie des faisceaux），Herman。
\item[{[8]}] J. Giraud：《\Term{descent}{下降}方法》（Méthode de la Descente），Mémoire de la S.M.F.
\item[{[9]}] J. Giraud：《非交换\Term{homology}{同调}\Term{algebra}{代数}》（Algèbre homologique non commutative），待出版。
\item[{[10]}] A. Grothendieck：《论\Term{algebraic-variety}{代数簇}的 de Rham \Term{cohomology}{上同调}》（On the De Rham Cohomology of Algebraic Varieties），I.H.E.S. 第 29 号。
\item[{[11]}] A. Grothendieck：《\Term{homology}{同调}\Term{algebra}{代数}的若干问题》（Sur quelques points d'Algèbre Homologique），Tohoku Math. Journal。
\item[{[12]}] R. Hartshorne：《\Term{residue}{留数}与\Term{duality}{对偶}》（Residues and Duality），Lecture Notes，第 20 号，Springer Verlag。
\item[{[13]}] L. Illusie：SGA 6 第 I 讲。
\item[{[14]}] S. Lubkin：《论 A. Weil 的一个猜想》（On a Conjecture of A. Weil），Am. J. of Math.，第 456 页，1967。
\item[{[15]}] G. Segal：《分类\Term{space}{空间}与\Term{spectral-sequence}{谱序列}》（Classifying spaces and Spectral Sequences），I.H.E.S. 第 34 号。
\item[{[16]}] Cartan 讨论班，1957--1957。
\item[{[17]}] D. Sullivan：《几何\Term{topology}{拓扑}，第一部分》（Geometric Topology, part I），油印讲义，M.I.T.，1970。
\end{enumerate}
\SourcePage{70}
